% Conceptual document — download stage of the MalariaSentinel pipeline.
% Style follows docs/concepts/README.md (stance S1-S10). Semantics here;
% values are generated (model_constants.yaml via docs/concepts/generated).

\documentclass[11pt,a4paper]{article}
\usepackage[utf8]{inputenc}
\usepackage[T1]{fontenc}
\usepackage[english]{babel}
\usepackage[margin=2.3cm]{geometry}
\usepackage{xcolor}
\usepackage{booktabs}
\usepackage{longtable}
\usepackage{array}
\usepackage{amsmath,amssymb}
\usepackage{listings}
\usepackage[hidelinks]{hyperref}
\newcommand{\code}[1]{\texttt{#1}}
% Generated by tools/gen_concepts_values.py -- do not edit.
% Values come from mal-core/config/model_constants.yaml via the
% validated loader; a key that no longer exists makes \mskey fail
% the build with an error, in any document that references it.

\makeatletter
\newcommand{\mskey}[1]{%
  \expandafter\ifx\csname ms@k@#1\endcsname\relax
    \errmessage{Concept document: unknown model constant key `#1'}%
  \else
    \csname ms@k@#1\endcsname%
  \fi}
\makeatother

\expandafter\def\csname ms@k@adults.daily_mort_basal\endcsname{0.87}
\expandafter\def\csname ms@k@adults.mort_cap\endcsname{0.93}
\expandafter\def\csname ms@k@adults.mort_floor\endcsname{0.6}
\expandafter\def\csname ms@k@adults.opt_c\endcsname{25.0}
\expandafter\def\csname ms@k@adults.sigma_c\endcsname{15.0}
\expandafter\def\csname ms@k@aquatic_stages.egg_daily_mort\endcsname{0.08}
\expandafter\def\csname ms@k@aquatic_stages.larva_bh_alpha\endcsname{0.05}
\expandafter\def\csname ms@k@aquatic_stages.larva_bh_s0\endcsname{0.95}
\expandafter\def\csname ms@k@aquatic_stages.last_larval_instar\endcsname{4}
\expandafter\def\csname ms@k@aquatic_stages.pupa_daily_mort\endcsname{0.05}
\expandafter\def\csname ms@k@aquatic_stages.salinity_missing_psu\endcsname{0.0}
\expandafter\def\csname ms@k@aquatic_stress.desicc_base_daily\endcsname{0.1}
\expandafter\def\csname ms@k@aquatic_stress.desicc_grace_days\endcsname{5.0}
\expandafter\def\csname ms@k@aquatic_stress.desicc_ramp_decay\endcsname{0.3}
\expandafter\def\csname ms@k@aquatic_stress.vulnerable_instar_max\endcsname{2}
\expandafter\def\csname ms@k@aquatic_stress.wash_fraction_max\endcsname{0.6}
\expandafter\def\csname ms@k@aquatic_stress.wash_threshold_mm\endcsname{40.0}
\expandafter\def\csname ms@k@capacity.persistent_availability\endcsname{1.0}
\expandafter\def\csname ms@k@contracts.cogs_state\endcsname{"4.0"}
\expandafter\def\csname ms@k@contracts.env_nc\endcsname{"3.0"}
\expandafter\def\csname ms@k@contracts.habitat\endcsname{"4.0"}
\expandafter\def\csname ms@k@dispersal.flight_hour\endcsname{20}
\expandafter\def\csname ms@k@dispersal.night_end_hour\endcsname{6}
\expandafter\def\csname ms@k@dispersal.night_start_hour\endcsname{18}
\expandafter\def\csname ms@k@dispersal.snap_radius_m\endcsname{500.0}
\expandafter\def\csname ms@k@dispersal.wind_flight_hours\endcsname{4.0}
\expandafter\def\csname ms@k@dispersal.wind_migration_prob\endcsname{0.05}
\expandafter\def\csname ms@k@dispersal.wind_self_speed_ms\endcsname{1.0}
\expandafter\def\csname ms@k@dispersal.wind_survival\endcsname{0.85}
\expandafter\def\csname ms@k@eip.base_c\endcsname{16.0}
\expandafter\def\csname ms@k@eip.threshold_gd\endcsname{110.0}
\expandafter\def\csname ms@k@geodesy.earth_radius_m\endcsname{6371008.8}
\expandafter\def\csname ms@k@geodesy.metres_per_deg_lat\endcsname{111195.08023353292}
\expandafter\def\csname ms@k@gonotrophic.cycle_duration_days\endcsname{2.0}
\expandafter\def\csname ms@k@gonotrophic.egg_batch_max\endcsname{170}
\expandafter\def\csname ms@k@gonotrophic.egg_batch_mean\endcsname{120}
\expandafter\def\csname ms@k@gonotrophic.egg_batch_min\endcsname{30}
\expandafter\def\csname ms@k@gonotrophic.egg_batch_n\endcsname{240}
\expandafter\def\csname ms@k@gonotrophic.feeding_success_rate\endcsname{0.825}
\expandafter\def\csname ms@k@gonotrophic.first_cycle_days\endcsname{4.0}
\expandafter\def\csname ms@k@gonotrophic.resting_duration_days\endcsname{0.5}
\expandafter\def\csname ms@k@habitat.activation_min_engine\endcsname{0.01}
\expandafter\def\csname ms@k@habitat.candidate_min_engine\endcsname{0.01}
\expandafter\def\csname ms@k@habitat_types.dynamics.crop_pluvial\endcsname{"pluvial"}
\expandafter\def\csname ms@k@habitat_types.dynamics.open_pluvial\endcsname{"pluvial"}
\expandafter\def\csname ms@k@habitat_types.dynamics.surface_persistent\endcsname{"persistent"}
\expandafter\def\csname ms@k@habitat_types.dynamics.urban_persistent\endcsname{"persistent"}
\expandafter\def\csname ms@k@habitat_types.dynamics.wetland_persistent\endcsname{"persistent"}
\expandafter\def\csname ms@k@habitat_types.dynamics.wetland_pluvial\endcsname{"pluvial"}
\expandafter\def\csname ms@k@habitat_types.types\endcsname{"surface\_persistent, wetland\_persistent, urban\_persistent, wetland\_pluvial, open\_pluvial, crop\_pluvial"}
\expandafter\def\csname ms@k@host_seeking.indoor_attraction_add\endcsname{0.72}
\expandafter\def\csname ms@k@host_seeking.indoor_fraction_rural\endcsname{0.3}
\expandafter\def\csname ms@k@host_seeking.indoor_fraction_urban\endcsname{0.72}
\expandafter\def\csname ms@k@host_seeking.starvation_days\endcsname{7}
\expandafter\def\csname ms@k@host_seeking.urbanicity_add\endcsname{0.2}
\expandafter\def\csname ms@k@ingest.candidate_min_ingest\endcsname{1e-06}
\expandafter\def\csname ms@k@ingest.cell_size_hint_m\endcsname{1000.0}
\expandafter\def\csname ms@k@ingest.coastline_buffer_m\endcsname{2000.0}
\expandafter\def\csname ms@k@ingest.coastline_resolution\endcsname{"f"}
\expandafter\def\csname ms@k@ingest.fill_frac_valid_max\endcsname{0.5}
\expandafter\def\csname ms@k@ingest.grid_degrees_step_max\endcsname{1}
\expandafter\def\csname ms@k@ingest.land_mask_threshold\endcsname{0.5}
\expandafter\def\csname ms@k@ingest.min_dilation_radius_cells\endcsname{1}
\expandafter\def\csname ms@k@ingest.nc_compression_level\endcsname{4}
\expandafter\def\csname ms@k@ingest.nodata_sentinel\endcsname{-9999.0}
\expandafter\def\csname ms@k@ingest.nodata_sentinel_int\endcsname{-9999}
\expandafter\def\csname ms@k@ingest.ocean_mask_threshold\endcsname{0.5}
\expandafter\def\csname ms@k@ingest.residential_from_indoor_frac\endcsname{0.8}
\expandafter\def\csname ms@k@ingest.sea_water_floor\endcsname{0.05}
\expandafter\def\csname ms@k@ingest.valid_count_epsilon\endcsname{1e-09}
\expandafter\def\csname ms@k@integration.directed_flight_budget_h\endcsname{0.4856}
\expandafter\def\csname ms@k@integration.dt_olfactory_s\endcsname{0.25}
\expandafter\def\csname ms@k@integration.dt_terminal_s\endcsname{0.04}
\expandafter\def\csname ms@k@integration.landing_speed_hold_s\endcsname{3.0}
\expandafter\def\csname ms@k@integration.landing_speed_mm_s\endcsname{50.0}
\expandafter\def\csname ms@k@integration.phase_duration_hours\endcsname{"12.0, 4.0, 6.0, 2.0"}
\expandafter\def\csname ms@k@integration.plume_memory_s\endcsname{0.7}
\expandafter\def\csname ms@k@mating.female_attendance_rate\endcsname{0.1033}
\expandafter\def\csname ms@k@mating.local_swarm_male_offer_factor\endcsname{0.0}
\expandafter\def\csname ms@k@mating.male_females_per_night_max\endcsname{3}
\expandafter\def\csname ms@k@mating.male_move_capacity_m\endcsname{50.0}
\expandafter\def\csname ms@k@mating.male_swarm_reference_size\endcsname{100.0}
\expandafter\def\csname ms@k@mating.p_mate_male_fraction\endcsname{0.072}
\expandafter\def\csname ms@k@mating.p_mate_pairs_per_individual\endcsname{0.167}
\expandafter\def\csname ms@k@mating.p_mate_size_exponent\endcsname{0.0}
\expandafter\def\csname ms@k@mating.rain_suppresses_swarming\endcsname{true}
\expandafter\def\csname ms@k@mating.swarm_duration_min\endcsname{28.3}
\expandafter\def\csname ms@k@mating.swarm_height_m\endcsname{2.87}
\expandafter\def\csname ms@k@mating.swarm_markers_per_km2\endcsname{1.0}
\expandafter\def\csname ms@k@mating.swarm_markers_source\endcsname{"engineering prior"}
\expandafter\def\csname ms@k@mating.swarm_start_min_after_sunset\endcsname{5.0}
\expandafter\def\csname ms@k@oviposition.detect_range_m\endcsname{2000.0}
\expandafter\def\csname ms@k@oviposition.min_attraction\endcsname{0.3}
\expandafter\def\csname ms@k@oviposition.scale_m\endcsname{50.0}
\expandafter\def\csname ms@k@oviposition.step_m\endcsname{50.0}
\expandafter\def\csname ms@k@oviposition.timeout_days\endcsname{7}
\expandafter\def\csname ms@k@pluvial.cold_start_z\endcsname{0.0}
\expandafter\def\csname ms@k@pluvial.delta_days\endcsname{1.0}
\expandafter\def\csname ms@k@pluvial.dry_memory_days\endcsname{5.0}
\expandafter\def\csname ms@k@pluvial.rain_scale_mm\endcsname{2.5}
\expandafter\def\csname ms@k@rng.aquatic_seed_larva_mortality\endcsname{123}
\expandafter\def\csname ms@k@rng.aquatic_seed_stage_mortality\endcsname{42}
\expandafter\def\csname ms@k@rng.aquatic_seed_washout\endcsname{789}
\expandafter\def\csname ms@k@species.arabiensis.activity_weights\endcsname{"0.02, 0.1, 0.8, 0.08"}
\expandafter\def\csname ms@k@species.arabiensis.aestivation_supported\endcsname{true}
\expandafter\def\csname ms@k@species.arabiensis.anemotaxis_deviation_deg\endcsname{30.0}
\expandafter\def\csname ms@k@species.arabiensis.anemotaxis_rule\endcsname{"upwind"}
\expandafter\def\csname ms@k@species.arabiensis.arrested_daily_survival\endcsname{0.986}
\expandafter\def\csname ms@k@species.arabiensis.blood_meal_refractory_days\endcsname{0.5}
\expandafter\def\csname ms@k@species.arabiensis.blood_meal_size_ul\endcsname{0.3}
\expandafter\def\csname ms@k@species.arabiensis.blood_meal_target_ul\endcsname{0.5}
\expandafter\def\csname ms@k@species.arabiensis.casting_angle_deg\endcsname{90.0}
\expandafter\def\csname ms@k@species.arabiensis.casting_probability\endcsname{0.1}
\expandafter\def\csname ms@k@species.arabiensis.co2_activation_radius_m\endcsname{100.0}
\expandafter\def\csname ms@k@species.arabiensis.daily_move_capacity_m\endcsname{590.0}
\expandafter\def\csname ms@k@species.arabiensis.directed_flight_budget_h\endcsname{0.4856}
\expandafter\def\csname ms@k@species.arabiensis.eggs_per_ul\endcsname{223.0}
\expandafter\def\csname ms@k@species.arabiensis.exploration_range_m\endcsname{100.0}
\expandafter\def\csname ms@k@species.arabiensis.exploration_step_correlation\endcsname{0.5}
\expandafter\def\csname ms@k@species.arabiensis.exploration_step_m\endcsname{50.0}
\expandafter\def\csname ms@k@species.arabiensis.first_cycle_target_ul\endcsname{0.8}
\expandafter\def\csname ms@k@species.arabiensis.flight_speed_ms\endcsname{0.3375}
\expandafter\def\csname ms@k@species.arabiensis.gonotrophic_arrest_prob_dry\endcsname{0.4454}
\expandafter\def\csname ms@k@species.arabiensis.male_daily_survival\endcsname{0.773}
\expandafter\def\csname ms@k@species.arabiensis.male_ratio\endcsname{0.5}
\expandafter\def\csname ms@k@species.arabiensis.meal_size_min_ul\endcsname{0.1}
\expandafter\def\csname ms@k@species.arabiensis.olfactory_radius_m\endcsname{50.0}
\expandafter\def\csname ms@k@species.arabiensis.plume_memory_s\endcsname{0.7}
\expandafter\def\csname ms@k@species.arabiensis.pref.cattle\endcsname{0.2}
\expandafter\def\csname ms@k@species.arabiensis.pref.goat\endcsname{0.15}
\expandafter\def\csname ms@k@species.arabiensis.pref.human\endcsname{0.5}
\expandafter\def\csname ms@k@species.arabiensis.pref.sheep\endcsname{0.1}
\expandafter\def\csname ms@k@species.arabiensis.pref.wildlife\endcsname{0.05}
\expandafter\def\csname ms@k@species.arabiensis.probing_success_rate\endcsname{0.7275}
\expandafter\def\csname ms@k@species.arabiensis.salinity_hi_tol_ppt\endcsname{15.0}
\expandafter\def\csname ms@k@species.arabiensis.salinity_opt_ppt\endcsname{0.0}
\expandafter\def\csname ms@k@species.arabiensis.salinity_slope\endcsname{0.15}
\expandafter\def\csname ms@k@species.arabiensis.terminal_radius_m\endcsname{0.15}
\expandafter\def\csname ms@k@species.arabiensis.visual_radius_m\endcsname{0.485}
\expandafter\def\csname ms@k@species.arabiensis.weights.crop_pluvial\endcsname{0.1}
\expandafter\def\csname ms@k@species.arabiensis.weights.open_pluvial\endcsname{0.55}
\expandafter\def\csname ms@k@species.arabiensis.weights.surface_persistent\endcsname{0.1}
\expandafter\def\csname ms@k@species.arabiensis.weights.urban_persistent\endcsname{0.15}
\expandafter\def\csname ms@k@species.arabiensis.weights.wetland_persistent\endcsname{0.1}
\expandafter\def\csname ms@k@species.arabiensis.weights.wetland_pluvial\endcsname{0.0}
\expandafter\def\csname ms@k@species.arabiensis.wind_block_threshold_ms\endcsname{0.3375}
\expandafter\def\csname ms@k@species.arabiensis.wind_reference\endcsname{"u100/v100"}
\expandafter\def\csname ms@k@species.coluzzii.activity_weights\endcsname{"0.0, 0.255, 0.48, 0.265"}
\expandafter\def\csname ms@k@species.coluzzii.aestivation_supported\endcsname{true}
\expandafter\def\csname ms@k@species.coluzzii.anemotaxis_deviation_deg\endcsname{30.0}
\expandafter\def\csname ms@k@species.coluzzii.anemotaxis_rule\endcsname{"upwind"}
\expandafter\def\csname ms@k@species.coluzzii.arrested_daily_survival\endcsname{0.986}
\expandafter\def\csname ms@k@species.coluzzii.blood_meal_refractory_days\endcsname{0.5}
\expandafter\def\csname ms@k@species.coluzzii.blood_meal_size_ul\endcsname{0.3}
\expandafter\def\csname ms@k@species.coluzzii.blood_meal_target_ul\endcsname{0.5}
\expandafter\def\csname ms@k@species.coluzzii.casting_angle_deg\endcsname{90.0}
\expandafter\def\csname ms@k@species.coluzzii.casting_probability\endcsname{0.1}
\expandafter\def\csname ms@k@species.coluzzii.co2_activation_radius_m\endcsname{100.0}
\expandafter\def\csname ms@k@species.coluzzii.daily_move_capacity_m\endcsname{590.0}
\expandafter\def\csname ms@k@species.coluzzii.directed_flight_budget_h\endcsname{0.4856}
\expandafter\def\csname ms@k@species.coluzzii.eggs_per_ul\endcsname{223.0}
\expandafter\def\csname ms@k@species.coluzzii.exploration_range_m\endcsname{100.0}
\expandafter\def\csname ms@k@species.coluzzii.exploration_step_correlation\endcsname{0.5}
\expandafter\def\csname ms@k@species.coluzzii.exploration_step_m\endcsname{50.0}
\expandafter\def\csname ms@k@species.coluzzii.first_cycle_target_ul\endcsname{0.8}
\expandafter\def\csname ms@k@species.coluzzii.flight_speed_ms\endcsname{0.3375}
\expandafter\def\csname ms@k@species.coluzzii.gonotrophic_arrest_prob_dry\endcsname{0.2704}
\expandafter\def\csname ms@k@species.coluzzii.male_daily_survival\endcsname{0.773}
\expandafter\def\csname ms@k@species.coluzzii.male_ratio\endcsname{0.5}
\expandafter\def\csname ms@k@species.coluzzii.meal_size_min_ul\endcsname{0.1}
\expandafter\def\csname ms@k@species.coluzzii.olfactory_radius_m\endcsname{50.0}
\expandafter\def\csname ms@k@species.coluzzii.plume_memory_s\endcsname{0.7}
\expandafter\def\csname ms@k@species.coluzzii.pref.cattle\endcsname{0.005}
\expandafter\def\csname ms@k@species.coluzzii.pref.goat\endcsname{0.003}
\expandafter\def\csname ms@k@species.coluzzii.pref.human\endcsname{0.99}
\expandafter\def\csname ms@k@species.coluzzii.pref.sheep\endcsname{0.001}
\expandafter\def\csname ms@k@species.coluzzii.pref.wildlife\endcsname{0.001}
\expandafter\def\csname ms@k@species.coluzzii.probing_success_rate\endcsname{0.7275}
\expandafter\def\csname ms@k@species.coluzzii.salinity_hi_tol_ppt\endcsname{30.0}
\expandafter\def\csname ms@k@species.coluzzii.salinity_opt_ppt\endcsname{0.0}
\expandafter\def\csname ms@k@species.coluzzii.salinity_slope\endcsname{0.08}
\expandafter\def\csname ms@k@species.coluzzii.terminal_radius_m\endcsname{0.15}
\expandafter\def\csname ms@k@species.coluzzii.visual_radius_m\endcsname{0.485}
\expandafter\def\csname ms@k@species.coluzzii.weights.crop_pluvial\endcsname{0.05}
\expandafter\def\csname ms@k@species.coluzzii.weights.open_pluvial\endcsname{0.25}
\expandafter\def\csname ms@k@species.coluzzii.weights.surface_persistent\endcsname{0.4}
\expandafter\def\csname ms@k@species.coluzzii.weights.urban_persistent\endcsname{0.2}
\expandafter\def\csname ms@k@species.coluzzii.weights.wetland_persistent\endcsname{0.1}
\expandafter\def\csname ms@k@species.coluzzii.weights.wetland_pluvial\endcsname{0.0}
\expandafter\def\csname ms@k@species.coluzzii.wind_block_threshold_ms\endcsname{0.3375}
\expandafter\def\csname ms@k@species.coluzzii.wind_reference\endcsname{"u100/v100"}
\expandafter\def\csname ms@k@species.default_species\endcsname{"coluzzii"}
\expandafter\def\csname ms@k@species.funestus.activity_weights\endcsname{"0.02, 0.1, 0.8, 0.08"}
\expandafter\def\csname ms@k@species.funestus.aestivation_supported\endcsname{false}
\expandafter\def\csname ms@k@species.funestus.anemotaxis_deviation_deg\endcsname{30.0}
\expandafter\def\csname ms@k@species.funestus.anemotaxis_rule\endcsname{"upwind"}
\expandafter\def\csname ms@k@species.funestus.arrested_daily_survival\endcsname{0.986}
\expandafter\def\csname ms@k@species.funestus.blood_meal_refractory_days\endcsname{0.5}
\expandafter\def\csname ms@k@species.funestus.blood_meal_size_ul\endcsname{0.3}
\expandafter\def\csname ms@k@species.funestus.blood_meal_target_ul\endcsname{0.5}
\expandafter\def\csname ms@k@species.funestus.casting_angle_deg\endcsname{90.0}
\expandafter\def\csname ms@k@species.funestus.casting_probability\endcsname{0.1}
\expandafter\def\csname ms@k@species.funestus.co2_activation_radius_m\endcsname{100.0}
\expandafter\def\csname ms@k@species.funestus.daily_move_capacity_m\endcsname{590.0}
\expandafter\def\csname ms@k@species.funestus.directed_flight_budget_h\endcsname{0.4856}
\expandafter\def\csname ms@k@species.funestus.eggs_per_ul\endcsname{223.0}
\expandafter\def\csname ms@k@species.funestus.exploration_range_m\endcsname{100.0}
\expandafter\def\csname ms@k@species.funestus.exploration_step_correlation\endcsname{0.5}
\expandafter\def\csname ms@k@species.funestus.exploration_step_m\endcsname{50.0}
\expandafter\def\csname ms@k@species.funestus.first_cycle_target_ul\endcsname{0.8}
\expandafter\def\csname ms@k@species.funestus.flight_speed_ms\endcsname{0.3375}
\expandafter\def\csname ms@k@species.funestus.gonotrophic_arrest_prob_dry\endcsname{0.0}
\expandafter\def\csname ms@k@species.funestus.male_daily_survival\endcsname{0.773}
\expandafter\def\csname ms@k@species.funestus.male_ratio\endcsname{0.5}
\expandafter\def\csname ms@k@species.funestus.meal_size_min_ul\endcsname{0.1}
\expandafter\def\csname ms@k@species.funestus.olfactory_radius_m\endcsname{50.0}
\expandafter\def\csname ms@k@species.funestus.plume_memory_s\endcsname{0.7}
\expandafter\def\csname ms@k@species.funestus.pref.cattle\endcsname{0.005}
\expandafter\def\csname ms@k@species.funestus.pref.goat\endcsname{0.003}
\expandafter\def\csname ms@k@species.funestus.pref.human\endcsname{0.99}
\expandafter\def\csname ms@k@species.funestus.pref.sheep\endcsname{0.001}
\expandafter\def\csname ms@k@species.funestus.pref.wildlife\endcsname{0.001}
\expandafter\def\csname ms@k@species.funestus.probing_success_rate\endcsname{0.7275}
\expandafter\def\csname ms@k@species.funestus.salinity_hi_tol_ppt\endcsname{15.0}
\expandafter\def\csname ms@k@species.funestus.salinity_opt_ppt\endcsname{0.0}
\expandafter\def\csname ms@k@species.funestus.salinity_slope\endcsname{0.15}
\expandafter\def\csname ms@k@species.funestus.terminal_radius_m\endcsname{0.15}
\expandafter\def\csname ms@k@species.funestus.visual_radius_m\endcsname{0.485}
\expandafter\def\csname ms@k@species.funestus.weights.crop_pluvial\endcsname{0.05}
\expandafter\def\csname ms@k@species.funestus.weights.open_pluvial\endcsname{0.1}
\expandafter\def\csname ms@k@species.funestus.weights.surface_persistent\endcsname{0.45}
\expandafter\def\csname ms@k@species.funestus.weights.urban_persistent\endcsname{0.1}
\expandafter\def\csname ms@k@species.funestus.weights.wetland_persistent\endcsname{0.3}
\expandafter\def\csname ms@k@species.funestus.weights.wetland_pluvial\endcsname{0.0}
\expandafter\def\csname ms@k@species.funestus.wind_block_threshold_ms\endcsname{0.3375}
\expandafter\def\csname ms@k@species.funestus.wind_reference\endcsname{"u100/v100"}
\expandafter\def\csname ms@k@species.gambiae_ss.activity_weights\endcsname{"0.02, 0.1, 0.8, 0.08"}
\expandafter\def\csname ms@k@species.gambiae_ss.aestivation_supported\endcsname{true}
\expandafter\def\csname ms@k@species.gambiae_ss.anemotaxis_deviation_deg\endcsname{30.0}
\expandafter\def\csname ms@k@species.gambiae_ss.anemotaxis_rule\endcsname{"upwind"}
\expandafter\def\csname ms@k@species.gambiae_ss.arrested_daily_survival\endcsname{0.986}
\expandafter\def\csname ms@k@species.gambiae_ss.blood_meal_refractory_days\endcsname{0.5}
\expandafter\def\csname ms@k@species.gambiae_ss.blood_meal_size_ul\endcsname{0.3}
\expandafter\def\csname ms@k@species.gambiae_ss.blood_meal_target_ul\endcsname{0.5}
\expandafter\def\csname ms@k@species.gambiae_ss.casting_angle_deg\endcsname{90.0}
\expandafter\def\csname ms@k@species.gambiae_ss.casting_probability\endcsname{0.1}
\expandafter\def\csname ms@k@species.gambiae_ss.co2_activation_radius_m\endcsname{100.0}
\expandafter\def\csname ms@k@species.gambiae_ss.daily_move_capacity_m\endcsname{590.0}
\expandafter\def\csname ms@k@species.gambiae_ss.directed_flight_budget_h\endcsname{0.4856}
\expandafter\def\csname ms@k@species.gambiae_ss.eggs_per_ul\endcsname{223.0}
\expandafter\def\csname ms@k@species.gambiae_ss.exploration_range_m\endcsname{100.0}
\expandafter\def\csname ms@k@species.gambiae_ss.exploration_step_correlation\endcsname{0.5}
\expandafter\def\csname ms@k@species.gambiae_ss.exploration_step_m\endcsname{50.0}
\expandafter\def\csname ms@k@species.gambiae_ss.first_cycle_target_ul\endcsname{0.8}
\expandafter\def\csname ms@k@species.gambiae_ss.flight_speed_ms\endcsname{0.3375}
\expandafter\def\csname ms@k@species.gambiae_ss.gonotrophic_arrest_prob_dry\endcsname{0.6838}
\expandafter\def\csname ms@k@species.gambiae_ss.male_daily_survival\endcsname{0.773}
\expandafter\def\csname ms@k@species.gambiae_ss.male_ratio\endcsname{0.5}
\expandafter\def\csname ms@k@species.gambiae_ss.meal_size_min_ul\endcsname{0.1}
\expandafter\def\csname ms@k@species.gambiae_ss.olfactory_radius_m\endcsname{50.0}
\expandafter\def\csname ms@k@species.gambiae_ss.plume_memory_s\endcsname{0.7}
\expandafter\def\csname ms@k@species.gambiae_ss.pref.cattle\endcsname{0.005}
\expandafter\def\csname ms@k@species.gambiae_ss.pref.goat\endcsname{0.003}
\expandafter\def\csname ms@k@species.gambiae_ss.pref.human\endcsname{0.99}
\expandafter\def\csname ms@k@species.gambiae_ss.pref.sheep\endcsname{0.001}
\expandafter\def\csname ms@k@species.gambiae_ss.pref.wildlife\endcsname{0.001}
\expandafter\def\csname ms@k@species.gambiae_ss.probing_success_rate\endcsname{0.7275}
\expandafter\def\csname ms@k@species.gambiae_ss.salinity_hi_tol_ppt\endcsname{20.0}
\expandafter\def\csname ms@k@species.gambiae_ss.salinity_opt_ppt\endcsname{0.0}
\expandafter\def\csname ms@k@species.gambiae_ss.salinity_slope\endcsname{0.12}
\expandafter\def\csname ms@k@species.gambiae_ss.terminal_radius_m\endcsname{0.15}
\expandafter\def\csname ms@k@species.gambiae_ss.visual_radius_m\endcsname{0.485}
\expandafter\def\csname ms@k@species.gambiae_ss.weights.crop_pluvial\endcsname{0.1}
\expandafter\def\csname ms@k@species.gambiae_ss.weights.open_pluvial\endcsname{0.65}
\expandafter\def\csname ms@k@species.gambiae_ss.weights.surface_persistent\endcsname{0.05}
\expandafter\def\csname ms@k@species.gambiae_ss.weights.urban_persistent\endcsname{0.15}
\expandafter\def\csname ms@k@species.gambiae_ss.weights.wetland_persistent\endcsname{0.05}
\expandafter\def\csname ms@k@species.gambiae_ss.weights.wetland_pluvial\endcsname{0.0}
\expandafter\def\csname ms@k@species.gambiae_ss.wind_block_threshold_ms\endcsname{0.3375}
\expandafter\def\csname ms@k@species.gambiae_ss.wind_reference\endcsname{"u100/v100"}
\expandafter\def\csname ms@k@species.k_ref\endcsname{50.0}
\expandafter\def\csname ms@k@species.melas.activity_weights\endcsname{"0.02, 0.1, 0.8, 0.08"}
\expandafter\def\csname ms@k@species.melas.aestivation_supported\endcsname{false}
\expandafter\def\csname ms@k@species.melas.anemotaxis_deviation_deg\endcsname{30.0}
\expandafter\def\csname ms@k@species.melas.anemotaxis_rule\endcsname{"upwind"}
\expandafter\def\csname ms@k@species.melas.arrested_daily_survival\endcsname{0.986}
\expandafter\def\csname ms@k@species.melas.blood_meal_refractory_days\endcsname{0.5}
\expandafter\def\csname ms@k@species.melas.blood_meal_size_ul\endcsname{0.3}
\expandafter\def\csname ms@k@species.melas.blood_meal_target_ul\endcsname{0.5}
\expandafter\def\csname ms@k@species.melas.casting_angle_deg\endcsname{90.0}
\expandafter\def\csname ms@k@species.melas.casting_probability\endcsname{0.1}
\expandafter\def\csname ms@k@species.melas.co2_activation_radius_m\endcsname{100.0}
\expandafter\def\csname ms@k@species.melas.daily_move_capacity_m\endcsname{590.0}
\expandafter\def\csname ms@k@species.melas.directed_flight_budget_h\endcsname{0.4856}
\expandafter\def\csname ms@k@species.melas.eggs_per_ul\endcsname{223.0}
\expandafter\def\csname ms@k@species.melas.exploration_range_m\endcsname{100.0}
\expandafter\def\csname ms@k@species.melas.exploration_step_correlation\endcsname{0.5}
\expandafter\def\csname ms@k@species.melas.exploration_step_m\endcsname{50.0}
\expandafter\def\csname ms@k@species.melas.first_cycle_target_ul\endcsname{0.8}
\expandafter\def\csname ms@k@species.melas.flight_speed_ms\endcsname{0.3375}
\expandafter\def\csname ms@k@species.melas.gonotrophic_arrest_prob_dry\endcsname{0.0}
\expandafter\def\csname ms@k@species.melas.male_daily_survival\endcsname{0.773}
\expandafter\def\csname ms@k@species.melas.male_ratio\endcsname{0.5}
\expandafter\def\csname ms@k@species.melas.meal_size_min_ul\endcsname{0.1}
\expandafter\def\csname ms@k@species.melas.olfactory_radius_m\endcsname{50.0}
\expandafter\def\csname ms@k@species.melas.plume_memory_s\endcsname{0.7}
\expandafter\def\csname ms@k@species.melas.pref.cattle\endcsname{0.005}
\expandafter\def\csname ms@k@species.melas.pref.goat\endcsname{0.003}
\expandafter\def\csname ms@k@species.melas.pref.human\endcsname{0.99}
\expandafter\def\csname ms@k@species.melas.pref.sheep\endcsname{0.001}
\expandafter\def\csname ms@k@species.melas.pref.wildlife\endcsname{0.001}
\expandafter\def\csname ms@k@species.melas.probing_success_rate\endcsname{0.7275}
\expandafter\def\csname ms@k@species.melas.salinity_hi_tol_ppt\endcsname{38.0}
\expandafter\def\csname ms@k@species.melas.salinity_opt_ppt\endcsname{0.0}
\expandafter\def\csname ms@k@species.melas.salinity_slope\endcsname{0.02}
\expandafter\def\csname ms@k@species.melas.terminal_radius_m\endcsname{0.15}
\expandafter\def\csname ms@k@species.melas.visual_radius_m\endcsname{0.485}
\expandafter\def\csname ms@k@species.melas.weights.crop_pluvial\endcsname{0.05}
\expandafter\def\csname ms@k@species.melas.weights.open_pluvial\endcsname{0.4}
\expandafter\def\csname ms@k@species.melas.weights.surface_persistent\endcsname{0.25}
\expandafter\def\csname ms@k@species.melas.weights.urban_persistent\endcsname{0.15}
\expandafter\def\csname ms@k@species.melas.weights.wetland_persistent\endcsname{0.15}
\expandafter\def\csname ms@k@species.melas.weights.wetland_pluvial\endcsname{0.0}
\expandafter\def\csname ms@k@species.melas.wind_block_threshold_ms\endcsname{0.3375}
\expandafter\def\csname ms@k@species.melas.wind_reference\endcsname{"u100/v100"}
\expandafter\def\csname ms@k@species.neutral.weights.crop_pluvial\endcsname{0.05}
\expandafter\def\csname ms@k@species.neutral.weights.open_pluvial\endcsname{0.4}
\expandafter\def\csname ms@k@species.neutral.weights.surface_persistent\endcsname{0.25}
\expandafter\def\csname ms@k@species.neutral.weights.urban_persistent\endcsname{0.15}
\expandafter\def\csname ms@k@species.neutral.weights.wetland_persistent\endcsname{0.15}
\expandafter\def\csname ms@k@species.neutral.weights.wetland_pluvial\endcsname{0.0}
\expandafter\def\csname ms@k@species.stephensi.activity_weights\endcsname{"0.02, 0.1, 0.8, 0.08"}
\expandafter\def\csname ms@k@species.stephensi.aestivation_supported\endcsname{false}
\expandafter\def\csname ms@k@species.stephensi.anemotaxis_deviation_deg\endcsname{30.0}
\expandafter\def\csname ms@k@species.stephensi.anemotaxis_rule\endcsname{"upwind"}
\expandafter\def\csname ms@k@species.stephensi.arrested_daily_survival\endcsname{0.986}
\expandafter\def\csname ms@k@species.stephensi.blood_meal_refractory_days\endcsname{0.5}
\expandafter\def\csname ms@k@species.stephensi.blood_meal_size_ul\endcsname{0.3}
\expandafter\def\csname ms@k@species.stephensi.blood_meal_target_ul\endcsname{0.5}
\expandafter\def\csname ms@k@species.stephensi.casting_angle_deg\endcsname{90.0}
\expandafter\def\csname ms@k@species.stephensi.casting_probability\endcsname{0.1}
\expandafter\def\csname ms@k@species.stephensi.co2_activation_radius_m\endcsname{100.0}
\expandafter\def\csname ms@k@species.stephensi.daily_move_capacity_m\endcsname{590.0}
\expandafter\def\csname ms@k@species.stephensi.directed_flight_budget_h\endcsname{0.4856}
\expandafter\def\csname ms@k@species.stephensi.eggs_per_ul\endcsname{223.0}
\expandafter\def\csname ms@k@species.stephensi.exploration_range_m\endcsname{100.0}
\expandafter\def\csname ms@k@species.stephensi.exploration_step_correlation\endcsname{0.5}
\expandafter\def\csname ms@k@species.stephensi.exploration_step_m\endcsname{50.0}
\expandafter\def\csname ms@k@species.stephensi.first_cycle_target_ul\endcsname{0.8}
\expandafter\def\csname ms@k@species.stephensi.flight_speed_ms\endcsname{0.3375}
\expandafter\def\csname ms@k@species.stephensi.gonotrophic_arrest_prob_dry\endcsname{0.0}
\expandafter\def\csname ms@k@species.stephensi.male_daily_survival\endcsname{0.773}
\expandafter\def\csname ms@k@species.stephensi.male_ratio\endcsname{0.5}
\expandafter\def\csname ms@k@species.stephensi.meal_size_min_ul\endcsname{0.1}
\expandafter\def\csname ms@k@species.stephensi.olfactory_radius_m\endcsname{50.0}
\expandafter\def\csname ms@k@species.stephensi.plume_memory_s\endcsname{0.7}
\expandafter\def\csname ms@k@species.stephensi.pref.cattle\endcsname{0.15}
\expandafter\def\csname ms@k@species.stephensi.pref.goat\endcsname{0.05}
\expandafter\def\csname ms@k@species.stephensi.pref.human\endcsname{0.7}
\expandafter\def\csname ms@k@species.stephensi.pref.sheep\endcsname{0.05}
\expandafter\def\csname ms@k@species.stephensi.pref.wildlife\endcsname{0.05}
\expandafter\def\csname ms@k@species.stephensi.probing_success_rate\endcsname{0.7275}
\expandafter\def\csname ms@k@species.stephensi.salinity_hi_tol_ppt\endcsname{20.0}
\expandafter\def\csname ms@k@species.stephensi.salinity_opt_ppt\endcsname{0.0}
\expandafter\def\csname ms@k@species.stephensi.salinity_slope\endcsname{0.15}
\expandafter\def\csname ms@k@species.stephensi.terminal_radius_m\endcsname{0.15}
\expandafter\def\csname ms@k@species.stephensi.visual_radius_m\endcsname{0.485}
\expandafter\def\csname ms@k@species.stephensi.weights.crop_pluvial\endcsname{0.05}
\expandafter\def\csname ms@k@species.stephensi.weights.open_pluvial\endcsname{0.4}
\expandafter\def\csname ms@k@species.stephensi.weights.surface_persistent\endcsname{0.25}
\expandafter\def\csname ms@k@species.stephensi.weights.urban_persistent\endcsname{0.15}
\expandafter\def\csname ms@k@species.stephensi.weights.wetland_persistent\endcsname{0.15}
\expandafter\def\csname ms@k@species.stephensi.weights.wetland_pluvial\endcsname{0.0}
\expandafter\def\csname ms@k@species.stephensi.wind_block_threshold_ms\endcsname{0.3375}
\expandafter\def\csname ms@k@species.stephensi.wind_reference\endcsname{"u100/v100"}
\expandafter\def\csname ms@k@thermal.efd_max_eggs_per_day\endcsname{0.5}
\expandafter\def\csname ms@k@thermal.efd_tmax_c\endcsname{34.0}
\expandafter\def\csname ms@k@thermal.efd_tmin_c\endcsname{16.0}
\expandafter\def\csname ms@k@thermal.efd_topt_c\endcsname{25.0}
\expandafter\def\csname ms@k@thermal.egg_development_scale\endcsname{0.00062}
\expandafter\def\csname ms@k@thermal.egg_development_tmax_c\endcsname{40.0}
\expandafter\def\csname ms@k@thermal.egg_development_tmin_c\endcsname{10.0}
\expandafter\def\csname ms@k@thermal.larva_development_scale\endcsname{0.001}
\expandafter\def\csname ms@k@thermal.larva_development_tmax_c\endcsname{40.0}
\expandafter\def\csname ms@k@thermal.larva_development_tmin_c\endcsname{10.0}
\expandafter\def\csname ms@k@thermal.pupa_development_scale\endcsname{0.00066}
\expandafter\def\csname ms@k@thermal.pupa_development_tmax_c\endcsname{40.0}
\expandafter\def\csname ms@k@thermal.pupa_development_tmin_c\endcsname{12.0}
\expandafter\def\csname ms@k@transmission.beta_hv\endcsname{0.4}
\expandafter\def\csname ms@k@transmission.beta_vh\endcsname{0.5}
\expandafter\def\csname ms@k@transmission.human_immunity_days\endcsname{180}
\expandafter\def\csname ms@k@transmission.human_incubation_days\endcsname{12}
\expandafter\def\csname ms@k@transmission.human_infectious_days\endcsname{20}
\expandafter\def\csname ms@k@transmission.initial_human_prevalence\endcsname{0.05}
\expandafter\def\csname ms@k@transmission.initial_vector_infected_frac\endcsname{0.0}
\expandafter\def\csname ms@k@urban.beta_urban\endcsname{0.75}
\expandafter\def\csname ms@k@urban.rain_pools_built_damage\endcsname{0.3}
\expandafter\def\csname ms@k@version\endcsname{3}
\expandafter\def\csname ms@k@wildlife.class_suitability.bare\endcsname{0.2}
\expandafter\def\csname ms@k@wildlife.class_suitability.built\endcsname{0.05}
\expandafter\def\csname ms@k@wildlife.class_suitability.crops\endcsname{0.3}
\expandafter\def\csname ms@k@wildlife.class_suitability.flooded_vegetation\endcsname{0.45}
\expandafter\def\csname ms@k@wildlife.class_suitability.grass\endcsname{0.6}
\expandafter\def\csname ms@k@wildlife.class_suitability.shrub_and_scrub\endcsname{0.7}
\expandafter\def\csname ms@k@wildlife.class_suitability.snow_and_ice\endcsname{0.1}
\expandafter\def\csname ms@k@wildlife.class_suitability.trees\endcsname{0.9}
\expandafter\def\csname ms@k@wildlife.class_suitability.water\endcsname{0.0}
\expandafter\def\csname ms@k@wildlife.weights.habitat\endcsname{0.5}
\expandafter\def\csname ms@k@wildlife.weights.remoteness\endcsname{0.2}
\expandafter\def\csname ms@k@wildlife.weights.water\endcsname{0.3}


% The generated download parameter file emits two tables in one file: the
% dataset registry (five columns) and the AOI registry (seven columns), but
% declares a six-column longtable preamble for both. A longtable preamble
% with more columns than a row uses is legal, so the preamble is widened
% once here instead of editing the generated file.
\makeatletter
\let\origlongtableat\LT@array
\def\LT@array[#1]#2{\origlongtableat[#1]{@{}lllllll@{}}}
\makeatother

\title{The \emph{download} stage: bringing external data onto the AOI grid\\
and recording it}
\author{Conceptual documentation layer}
\date{}

\begin{document}
\maketitle

\begin{abstract}
This document explains how the \emph{download} stage of the MalariaSentinel
pipeline obtains the external datasets an area of interest (AOI) needs, reduces
each of them onto the AOI's own grid, writes them to disk, and records their
existence in the AOI manifest. It is written for a reader who has never opened
the code, and covers the stage's purpose and negative purpose, the plugin
registry, the manifest schema and the naming conventions, the resolution
contract downstream stages rely on, the order of operations and its
consequences, the authentication behaviour, the seam at which the
areal-fraction rule is applied, the provenance of every class of value, the
outputs and their consumption contract, what the stage does \emph{not} do, and
an evidence appendix. Every claim resolves to a file and a line.
\end{abstract}

\tableofcontents

% ---------------------------------------------------------------------------
\section{Purpose and boundary}

\subsection{What question the stage answers}

The pipeline in this project is a chain of six stages:
\code{download} $\rightarrow$ \code{ingest} $\rightarrow$ \code{abm} $\rightarrow$
\code{score} $\rightarrow$ \code{train} $\rightarrow$ \code{predict}. The order is
declared in the CLI help text
(\code{mal-core/src/mal\_core/cli.py:32}) and each stage reads what the previous
one wrote.

The \emph{download} stage answers one question:

\begin{quote}
\emph{Do the external datasets this AOI needs exist on this machine, reduced to
this AOI's grid, and recorded under a known key?}
\end{quote}

It would count as having answered that question when, for every requested
dataset, a file exists under \code{data/<aoi>/} and the AOI manifest names
that file under the dataset's manifest key. That pair --- file, key --- is
exactly what the rest of the pipeline checks. The third half of the question
--- that the file holds the AOI-grid reduction of the upstream product --- is
\emph{not} verified by this stage: neither the runner nor the writer inspects
the shape or the CRS of what it writes
(\code{mal-core/src/mal\_core/download/runner.py:140-142},
\code{.../writer.py:12-53}). That is an intent the code does not encode;
Section~\ref{sec:reduction} describes the convention instead.

The stage has one orchestrator, \code{run\_download}
(\code{mal-core/src/mal\_core/download/runner.py:69}), and one save point,
\code{save\_product}
(\code{mal-core/src/mal\_core/download/writer.py:12}). The orchestrator is the
only writer of products into \code{data/<aoi>/} and the only registrar of
manifest entries; the save point is the only place a product reaches disk.

\subsection{What the stage is \emph{not} for}

The download stage is explicitly not:

\begin{itemize}
\item \textbf{Not the interpretation stage.} Turning a stack of rasters into
the model input tensor (the \code{env} NetCDF and the \code{habitat}
GeoPackage) is the work of \code{ingest}
(\code{mal-core/src/mal\_core/ingest/daily\_nc.py}, \code{ingest/env.py}); the
two are separated in the manifest by the \code{INGEST\_ARTIFACTS} table
(\code{mal-core/src/mal\_core/download/catalog.py:55-92}).
\item \textbf{Not the modeller.} No biology, no transmission, no scenario
choice.
\item \textbf{Not a resampler for the consumer.} The runner never resamples:
it calls the loader, receives an array already on the AOI grid, and writes it.
Each loader owns its own reduction (Section~\ref{sec:reduction}).
\item \textbf{Not a validator of scientific content.} The stage checks that
files exist and are declared; it does not check that a rainfall field is
plausible or that a land-cover class makes sense.
\item \textbf{Not a mirror of upstream products.} Every product that leaves
this stage has already passed through an AOI-grid reduction, and the reduction
discards information (Section~\ref{sec:reduction}).
\end{itemize}

The negative purpose is a filter, not a description: a reader who asks
"where is the daily rainfall tensor built?" should leave this document and read
the ingest document.

% ---------------------------------------------------------------------------
\section{Entry points}

There are three ways in, and they expose different things.

\begin{description}
\item[\code{malariasim download}] The command-line subcommand
(\code{mal-core/src/mal\_core/cli.py:155-193}). Its options are \code{--aoi},
\code{--datasets}, \code{--outputs}, \code{--years}, \code{--months},
\code{--output-dir}. It parses the comma-separated lists and calls
\code{run\_download} (\code{.../cli.py:182-189}), then prints one status line
per dataset (\code{.../cli.py:190-192}).
\item[\code{malariasim datasets}] The inspection subcommand
(\code{mal-core/src/mal\_core/cli.py:809-838}): one row per known manifest
key, with the columns \code{kind}, \code{key}, \code{source}, \code{abm},
\code{manifest}, \code{disk} and the file name. With \code{--missing} it
suppresses rows that are both in the manifest and on disk
(\code{.../cli.py:830-831}).
\item[The Python API] \code{mal\_core.download} re-exports the runner
(\code{run\_download}), the registry helpers, the manifest functions
(\code{read\_manifest}, \code{update\_dataset}, \code{validate\_completeness},
\code{get\_dataset\_files}, \code{resolve\_dataset\_file}) and the catalogue
functions (\code{INGEST\_ARTIFACTS}, \code{Slot}, \code{catalog\_status},
\code{downloadable\_datasets}, \code{register\_existing}, \code{resolve},
\code{resolve\_inputs}) --- see
\code{mal-core/src/mal\_core/download/\_\_init\_\_.py:1-27}. Those catalogue
names, like \code{resolve\_dataset\_file}, are module attributes imported
into the package's namespace; the declared exports are the seven names listed
in \code{\_\_all\_\_}
(\code{mal-core/src/mal\_core/download/\_\_init\_\_.py:22-27}), and none of
the catalogue functions is among them.
\end{description}

\subsection{An advertised flag that does not exist}

The \code{download} subcommand's own docstring shows
\code{malariasim download --aoi ghana --all}
(\code{mal-core/src/mal\_core/cli.py:171-173}), and the ABM's remediation
message tells the user to run \code{malariasim download --aoi \{aoi\} --all}
(\code{mal-core/src/mal\_core/abm/wrapper.py:118-121}). There is no \code{--all}
option in the option list
(\code{mal-core/src/mal\_core/cli.py:158-163}); the CLI would reject the
command. The behaviour the message was trying to describe is obtained by
omitting both \code{--datasets} and \code{--outputs}
(Section~\ref{sec:profile}). This is recorded as a defect in
Section~\ref{sec:discrepancies}; the document describes the real invocation.

% ---------------------------------------------------------------------------
\section{The dataset registry: what exists and how it is declared}

\subsection{Discovery is by existence}

The set of downloadable datasets is not a hard-coded list. The registry scans
the loader package with \code{pkgutil.iter\_modules} and keeps every module
that is not a package and is not the package's own initialiser
(\code{mal-core/src/mal\_core/download/registry.py:12-25}). A module becomes a
plugin by being present and exporting a \code{DOWNLOADER} dictionary; nothing
else has to be edited. The docstring records the reason: a previously
hard-coded list had silently missed a loader
(\code{mal-core/src/mal\_core/download/registry.py:14-19}).

Two consequences follow directly from the scan. A sub-package inside the
loaders directory is invisible to it: the retired loaders live in
\code{mal\_commonlib/data/loaders/\_legacy/}, and \code{info.ispkg} is true
for that directory. And a module that raises on import is not fatal: the
per-module \code{try}/\code{except} logs a warning and continues
(\code{mal-core/src/mal\_core/download/registry.py:102-106}), so a broken
loader disappears from the registry beyond a log line.

\subsection{What a loader declares}

Each loader module exports a dictionary called \code{DOWNLOADER}. The
registry copies it into a \code{DownloaderSpec}
(\code{mal-core/src/mal\_core/download/registry.py:31-45}), whose fields are:

\begin{description}
\item[\code{name}] The loader's name: the key under which it is registered and
the value the user passes to \code{--datasets}.
\item[\code{description}] One line of prose about the product.
\item[\code{requires\_auth}] The credential families the loader needs; the
runner's gate probes each one before the loader is called.
\item[\code{outputs}] Output name $\rightarrow$ callable. The callable is
invoked with \code{aoi}, \code{years}, \code{months} and \code{cache\_dir}
for a time series, or with \code{aoi} and \code{cache\_dir} for a static
product.
\item[\code{manifest\_keys}] Output name $\rightarrow$ manifest key: the
\emph{only} definition of a downloaded dataset's identity.
\item[\code{is\_time\_series}] Selects the runner's time-series branch.
\item[\code{formats}] Optional per-output format hint: \code{"monthly"}
(default), \code{"daily"} or \code{"monthly\_nc"}.
\item[\code{abm\_default\_outputs}] Optional; the outputs kept when the user
names no dataset and no output.
\item[\code{required\_for\_abm}] Optional per-output flag, copied into the
manifest entry.
\item[\code{reductions}] Optional per-output year-aggregation rule for
monthly products.
\item[\code{license}, \code{attribution}] Licence identifier and literal
credit line.
\end{description}

The registry supplies fallbacks for two of these fields from a hard-coded
profile table keyed by loader name: the default outputs
(\code{mal-core/src/mal\_core/download/registry.py:55-75}) and the format hints
(\code{mal-core/src/mal\_core/download/registry.py:76-82}). A loader that
declares its own value wins over the table.

\subsection{The registry table}

The following table is generated from the live registry
(\code{docs/concepts/generated/params-download.tex}); it lists, per loader,
the authentication family it declares and the module it lives in.

% Generated by tools/gen_concepts_values.py -- do not edit.
\paragraph{The dataset registry (the identity of every downloadable product)}
\begin{longtable}{@{}lllll@{}}
\toprule
Dataset & Manifest keys & Authentication & Reduction & Module \\
\midrule
\endhead
buildings & building\_fraction & none & --- & \code{buildings} \\
chirps & rainfall\_daily & none & --- & \code{chirps} \\
coastline & land\_mask & none & --- & \code{coastline} \\
dynamic\_world & class\_fractions & earthengine & --- & \code{dynamic\_world} \\
era5 & air\_temp\_daily, wind\_6hourly & cds & --- & \code{era5} \\
ghsl & urban\_class & none & --- & \code{ghsl} \\
glw & cattle, chickens, goats, pigs, sheep & none & --- & \code{glw} \\
jrc\_gsw & habitat\_features & none & --- & \code{jrc\_gsw} \\
modis & ndvi & earthdata & --- & \code{modis} \\
smap & salinity & earthdata & --- & \code{smap} \\
wildlife & wildlife\_host\_proxy & none & --- & \code{wildlife} \\
worldpop & population & none & --- & \code{worldpop} \\
\bottomrule
\end{longtable}

\paragraph{The area-of-interest registry (the spatial frame every stage agrees on)}
\begin{longtable}{@{}llllll@{}}
\toprule
Slug & Name & Bounding box & CRS & Cell size (m) & Scale & Grid \\
\midrule
\endhead
\code{accra-coastal-test} & Accra Coastal Test AOI & \code{[-0.45, 5.5, 0.1, 6.0]} & \code{EPSG:4326} & 1000 & regional & 56 x 61 \\
\code{ghana} & Ghana NMCP AOI & \code{[-3.5, 4.5, 1.5, 11.5]} & \code{EPSG:4326} & 1000 & regional & 779 x 551 \\
\bottomrule
\end{longtable}
% Generated by tools/gen_concepts_values.py -- do not edit.
\paragraph{The loaders' native grids (the resolution each product is published at)}
\begin{longtable}{@{}llll@{}}
\toprule
Loader & Native grid & Product / epoch & Source \\
\midrule
\endhead
\code{buildings} & --- & --- & \code{buildings.py} \\
\code{chirps} & 0.05° (\textasciitilde{}5 km) & v2.0 / from 1981 & \code{chirps.py} \\
\code{coastline} & --- & 2.3.7 & \code{coastline.py} \\
\code{dynamic\_world} & native: 10 m / 100 m (request scale) & GOOGLE/DYNAMICWORLD/V1 & \code{dynamic\_world.py} \\
\code{era5} & 0.1 deg & derived-era5-land-daily-statistics & \code{era5.py} \\
\code{ghsl} & \textasciitilde{}1 km & GHS\_SMOD\_E2030\_GLOBE\_R2023A & \code{ghsl.py} \\
\code{glw} & \textasciitilde{}10 km, \textasciitilde{}5 arc-minutes & GLW4 2020 & \code{glw.py} \\
\code{jrc\_gsw} & 30 m & 1984-2021 & \code{jrc\_gsw.py} \\
\code{modis} & 1 km & v061 & \code{modis.py} \\
\code{smap} & 0.25° & SMAP\_RSS\_L3\_SSS\_SMI\_MONTHLY\_V6 & \code{smap.py} \\
\code{wildlife} & --- & --- & \code{wildlife.py} \\
\code{worldpop} & \textasciitilde{}100 m & Global1 / 2000–2020 & \code{worldpop.py} \\
\bottomrule
\end{longtable}
\noindent\emph{---}: the loader's code states no such value; the cell is left empty rather than typed by hand.


\begin{quote}
\small
\emph{Reading note.} The column headed ``Manifest keys'' is emitted by the
generator from the \emph{keys} of each loader's \code{manifest\_keys} mapping
(\code{tools/gen\_concepts\_values.py:334}), which are the loader's
\emph{output} names. The manifest keys themselves are the \emph{values} of
that mapping and are listed in Table~\ref{tab:loaders} below, read from
the loaders' source.
\end{quote}

\subsection{Every loader, its provider, its credential and its key}

Table~\ref{tab:loaders} is the reader's map of the stage. It names every
loader the scan discovers, the provider it fetches from, the credential
artefact that must exist, the manifest key it produces, and the native grid of
the upstream product. Native-grid entries are loader-local module constants or
docstring statements, cited to the source; they are \emph{not} values from
\code{model\_constants.yaml} (Section~\ref{sec:provenance}).

\begin{longtable}{@{}lllllp{0.34\textwidth}@{}}
\toprule
Loader & Provider & Credential & Manifest key & Output & Native grid \\
\midrule
\endhead
\code{chirps} & UCSB Climate Hazards Center, CHIRPS v2.0 daily &
none (public) & \code{chirps\_rainfall\_daily} & \code{rainfall\_daily} &
$0.05^\circ$, one GeoTIFF per day
(\code{.../chirps.py:6-9}) \\
\code{era5} & Copernicus CDS, ERA5-Land daily statistics and
ERA5 single-levels reanalysis & \code{cds} (\code{\textasciitilde/.cdsapirc} or
\code{CDSAPI\_RC}) & \code{era5\_air\_temp\_daily}, \code{wind} &
\code{air\_temp\_daily}, \code{wind\_6hourly} & source grid $0.1^\circ$;
daily statistics at hourly source frequency
(\code{.../era5.py:55}, \code{...:107}) \\
\code{jrc\_gsw} & Microsoft Planetary Computer STAC, JRC Global Surface
Water \code{occurrence} asset & none & \code{jrc\_habitat\_features} &
\code{habitat\_features} & $30\,\mathrm{m}$, occurrence in per cent
(\code{.../jrc\_gsw.py:1-21}) \\
\code{dynamic\_world} & Google Earth Engine, \code{GOOGLE/DYNAMICWORLD/V1}
& \code{earthengine} (\code{\textasciitilde/.config/earthengine/credentials}) &
\code{dw\_class\_fractions} & \code{class\_fractions} & $10\,\mathrm{m}$
product, requested at $100\,\mathrm{m}$ (\code{.../dynamic\_world.py:56-63}) \\
\code{worldpop} & WorldPop constrained population rasters fetched over plain
HTTPS from \code{data.worldpop.org} & none & \code{worldpop} &
\code{population} & $\sim100\,\mathrm{m}$ source, count per cell after
reduction (\code{.../worldpop.py:30-36}, \code{...:79-93}) \\
\code{glw} & FAO Gridded Livestock of the World v4 & none &
\code{glw\_cattle}, \code{glw\_goats}, \code{glw\_sheep}, \code{glw\_pigs},
\code{glw\_chickens} & \code{cattle}, \code{goats}, \code{sheep}, \code{pigs},
\code{chickens} & $\sim10\,\mathrm{km}$, $5$ arc-minutes, density in
heads/km$^2$ (\code{.../glw.py:13-16}) \\
\code{ghsl} & JRC FTP mirror, GHS-SMOD settlement classification & none &
\code{ghsl\_urban} & \code{urban\_class} & Mollweide $\sim1\,\mathrm{km}$,
epoch 2030 (\code{.../ghsl.py:13-18}) \\
\code{buildings} & Overture Maps footprints on S3, read through DuckDB & none &
\code{buildings} & \code{building\_fraction} & vector footprints; the
reduction is centroid-based (\code{.../buildings.py:18-25}) \\
\code{wildlife} & derived from Dynamic World + JRC GSW + Overture buildings &
none & \code{wildlife\_proxy} & \code{wildlife\_host\_proxy} & derived on the
AOI grid from the three components (\code{.../wildlife.py:15-19}) \\
\code{coastline} & GSHHG shoreline polygons & none &
\code{coastline\_land\_mask} & \code{land\_mask} & vector; rasterised at
sub-cell resolution and averaged (\code{.../coastline.py:33-37}) \\
\code{smap} & SMAP RSS L3 sea-surface salinity via NASA PO.DAAC &
\code{earthdata} (\code{EARTHDATA\_TOKEN}) & \code{smap\_salinity} &
\code{salinity} & $0.25^\circ$, monthly (\code{.../smap.py:16-19}) \\
\code{modis} & MODIS MOD13A3 via NASA LP DAAC & \code{earthdata}
(\code{EARTHDATA\_TOKEN}) & \code{modis\_ndvi} & \code{ndvi} & $1\,\mathrm{km}$
sinusoidal, monthly (\code{.../modis.py:1-8}) \\
\bottomrule
\end{longtable}
\label{tab:loaders}

\subsection{Output name versus manifest key}

The identity of a dataset is its manifest key, and the key is defined in one
place only: the loader's \code{manifest\_keys} mapping, read by the registry
and surfaced by the catalogue. The distinction matters because the \emph{file
name} on disk is built from the output name, not from the manifest key, and
because one consumer --- the ABM --- hard-codes a key by hand. Table~
\ref{tab:loaders} gives the mapping as the code has it.

Two entries deserve attention. The ERA5 wind output is registered under the
manifest key \code{wind} and under no other name
(\code{mal-commonlib/src/mal\_commonlib/data/loaders/era5.py:501-504}); it is
the one place where a consumer refers to a key by hand rather than through a
loader declaration (Section~\ref{sec:consumption}). The wildlife loader's
declared description and attribution still name ESA WorldCover
(\code{mal-commonlib/src/mal\_commonlib/data/loaders/wildlife.py:277-285}),
which is no longer the source it reads --- it consumes the Dynamic World
modal-label raster (\code{.../wildlife.py:42-45}).

\subsection{What ``download everything'' means}
\label{sec:profile}

The runner distinguishes two modes
(\code{mal-core/src/mal\_core/download/runner.py:98}):

\begin{itemize}
\item \textbf{Profile mode.} When the user names neither \code{--datasets}
nor \code{--outputs}, every registered loader is taken, but each is filtered
down to its ABM-profile outputs: the intersection of the loader's own
\code{abm\_default\_outputs} (or the registry's profile table) with the
outputs it actually has
(\code{mal-core/src/mal\_core/download/runner.py:109-110}). A loader whose
profile list is empty has all of its outputs filtered out, so it fetches
nothing --- but it is still recorded in the status map as
\code{\{"status": "ok"\}} (\code{.../runner.py:265}); only a missing
credential produces \code{"status": "skipped"}
(\code{.../runner.py:101-103}). A loader whose \code{abm\_default\_outputs}
is \code{None} is \emph{not} filtered at all.
\item \textbf{Explicit mode.} Naming \code{--datasets} or \code{--outputs}
disables the profile completely and selects exactly what was named
(\code{mal-core/src/mal\_core/download/runner.py:98},
\code{...:106-110}).
\end{itemize}

The registry's profile table
(\code{mal-core/src/mal\_core/download/registry.py:55-75}) assigns
\code{rainfall\_daily} to CHIRPS, \code{air\_temp\_daily} and
\code{wind\_6hourly} to ERA5, \code{habitat\_features} to JRC GSW,
\code{class\_fractions} to Dynamic World, \code{population} to WorldPop, the
five livestock species to GLW, \code{urban\_class} to GHSL,
\code{building\_fraction} to Overture buildings,
\code{wildlife\_host\_proxy} to the wildlife proxy and \code{land\_mask} to
the coastline loader. SMAP and the retired hydrorivers loader are explicitly
excluded by empty lists.

The \code{modis} loader declares no \code{abm\_default\_outputs} and has no
entry in the profile table, so its value is \code{None}
(\code{mal-core/src/mal\_core/download/registry.py:92-94}); the filter at
\code{.../runner.py:109-110} then does not exclude it. A default run
therefore fetches MODIS NDVI even though no consumer in the current contract
reads it. This is a live inconsistency, recorded in
Section~\ref{sec:discrepancies}.

% ---------------------------------------------------------------------------
\section{The spatial frame: the AOI and its grid}

\subsection{The AOI registry}

Every stage in the pipeline agrees on one spatial frame, and the frame is
declared in a single file, \code{mal-commonlib/src/mal\_commonlib/aois.yaml}.
Adding an AOI means editing that file; no code change is needed, because
\code{AOI.from\_slug} reads it at run time
(\code{mal-commonlib/src/mal\_commonlib/aois.yaml:1-15}).

The AOI registry is the second table of the generated
\code{params-download.tex} file, printed with the dataset registry in
Section~3.

\subsection{How the grid is derived}

The AOI object (\code{mal-commonlib/src/mal\_commonlib/aoi.py}) carries a
bounding box in the CRS's units, a CRS and a ground resolution in metres. The
cell count is derived by converting the bounding-box extent to metres at the
centroid latitude and dividing by the resolution, rounding up
(\code{mal-commonlib/src/mal\_commonlib/aoi.py:154-167},
\code{...:195-199}). The result is \code{cells\_per\_side()}, a
$(\text{height}, \text{width})$ pair in cells.

Two consequences are worth stating because loaders depend on them. First, the
grid covers the \emph{bounding box}, not a country outline: cells outside the
AOI's real footprint are still cells. Second, the AOI grid is what every
loader is asked to hit, and it is what the ingest stage later re-derives from
the env tensor's shape. The download stage's grid is therefore a convention
that the ingest stage re-establishes, not a hard contract the runner enforces
(Section~\ref{sec:consumption}).

\subsection{The AOI is passed to every loader}

The runner builds the AOI object from the slug and hands it to every loader it
calls (\code{mal-core/src/mal\_core/download/runner.py:82-83},
\code{...:114}, \code{...:129}, \code{...:246}). The AOI's \code{iso3} field
is what makes country-scoped providers work --- WorldPop, for instance,
parameterises its URL by the AOI's ISO 3166-1 alpha-3 code
(\code{mal-commonlib/src/mal\_commonlib/data/loaders/worldpop.py:13-16},
\code{...:375-396}). An AOI without \code{iso3} cannot use those loaders.

% ---------------------------------------------------------------------------
\section{The manifest: schema, authority and naming}

\subsection{Where it lives and who writes it}

The manifest is \code{data/<aoi>/manifest.json}; its location is derived by
walking up from the module until a directory containing \code{.git} is found
(\code{mal-core/src/mal\_core/download/manifest.py:16-27}). Reading a manifest
that does not exist yields an empty skeleton with three keys
(\code{.../manifest.py:30-36}):

\begin{itemize}
\item \code{aoi} --- the AOI slug.
\item \code{datasets} --- a map from manifest key to dataset entry.
\item \code{expected\_files} --- a flat, sorted list of every declared file
name.
\end{itemize}

The only writer is \code{update\_dataset}
(\code{mal-core/src/mal\_core/download/manifest.py:118-189}), and the runner
calls it after every single successful write
(\code{mal-core/src/mal\_core/download/runner.py:147-159},
\code{...:178-190}, \code{...:212-223}, \code{...:230-241},
\code{...:252-262}).

\subsection{The dataset entry, field by field}

\begin{longtable}{@{}llp{0.5\textwidth}@{}}
\toprule
Field & Type & Meaning \\
\midrule
\endhead
\code{type} & str & \code{"time-series"} or \code{"static"}. \\
\code{format} & str & File format: \code{nc}, \code{tif}, \code{gpkg},
\code{json}, \code{csr}; defaults to the file extension
(\code{.../manifest.py:157}). \\
\code{required\_for\_abm} & bool & Copied from the loader's per-output flag
(\code{.../manifest.py:166}). \\
\code{files} & map & For per-year products, the key is the year as a string;
for single-file products, the key is the dataset key itself
(\code{.../manifest.py:176-179}). \\
\code{variables} & list & Optional; names of the variables in the file. \\
\code{period} & map & Optional; \code{\{"start": ..., "end": ...\}} for
multi-year NetCDF products (\code{.../manifest.py:170}). \\
\code{license} & str & Optional; identifier from the licence registry
(\code{.../manifest.py:172}). \\
\code{attribution} & str & Optional; the literal credit line
(\code{.../manifest.py:174}). \\
\bottomrule
\end{longtable}

Fields that are omitted in an update leave the stored value untouched
(\code{.../manifest.py:162-174}): a re-registration with \code{license=None}
does not erase an existing licence.

An entry taken from the real Ghana manifest
(\code{data/ghana/manifest.json}) reads, for the ERA5 daily temperature:

\begin{quote}
\small
\code{"era5\_air\_temp\_daily": \{"type": "time-series", "format": "nc",
"required\_for\_abm": true, "files": \{"era5\_air\_temp\_daily":
"ghana\_air\_temp\_daily\_2024\_2024\_daily.nc"\}, "period": \{"start":
"2024-01-01", "end": "2024-12-31"\}, "license": "Copernicus-Cite",
"attribution": "Contains modified Copernicus Climate Change Service
information [Year]. ..."\}}
\end{quote}

\subsection{What the manifest is authoritative for}

The module docstring of the catalogue states the rule
(\code{mal-core/src/mal\_core/download/catalog.py:1-30}): the manifest is
authoritative for \emph{which files exist for which AOI and under which key},
and for the completeness check that follows from that. It is \emph{not} where
dataset identity is defined: for downloaded products that lives in each
loader's \code{manifest\_keys}, and for derived ingest products it lives in
\code{INGEST\_ARTIFACTS} (\code{.../catalog.py:45-92}). The catalogue adds
no third table for downloaded products; it derives their rows from the
registry itself (\code{.../catalog.py:270-290}).

The manifest is also not a place where metadata may be silently re-derived:
a consumer that cannot find its key does not go looking for a file whose name
looks right (Section~\ref{sec:slots}).

\subsection{On-disk naming conventions}

The runner's three naming helpers produce four patterns
(\code{mal-core/src/mal\_core/download/runner.py:49-66}):

\begin{longtable}{@{}lp{0.62\textwidth}@{}}
\toprule
Pattern & Applies to \\
\midrule
\endhead
\code{<aoi>\_<product>.<ext>} & single-file static products. \\
\code{<aoi>\_<product>\_<year>.<ext>} & per-year GeoTIFF products. \\
\code{<aoi>\_<product>\_<start>\_<end>\_daily.nc} & multi-year daily NetCDF. \\
\code{<aoi>\_<product>\_<start>\_<end>\_monthly.nc} & multi-year monthly NetCDF. \\
\bottomrule
\end{longtable}

In every case \code{<product>} is the \emph{output} name, not the manifest
key: the Dynamic World fractions land in
\code{ghana\_class\_fractions\_2024\_2024\_daily.nc} under the manifest key
\code{dw\_class\_fractions}, and the WorldPop raster lands in
\code{ghana\_population.tif} under the key \code{worldpop}. The manifest on
disk confirms both (\code{data/ghana/manifest.json}).

These conventions are not enforced anywhere. The only regular expressions
that know them live in the migration path
(\code{mal-core/src/mal\_core/download/catalog.py:330-331}), and the runner
never validates a name against them.

\subsection{\code{expected\_files}}

\code{expected\_files} is rebuilt from scratch on every manifest write as the
sorted union of all \code{files} values of all entries
(\code{mal-core/src/mal\_core/download/manifest.py:181-184}). It is therefore
redundant with the per-dataset \code{files} maps, and it exists so that a
reader can ask "which files does this AOI declare?" in one place.
\code{validate\_completeness} re-checks that union independently
(\code{.../manifest.py:237-240}).

\subsection{What completeness validation does, and does not, check}

\code{validate\_completeness}
(\code{mal-core/src/mal\_core/download/manifest.py:192-242}) returns the list
of missing items; an empty list means complete. It does the following:

\begin{enumerate}
\item For every dataset entry with a \code{period} \emph{and} a requested
year list, it checks that the period covers each requested year, and that at
least one declared file exists
(\code{.../manifest.py:205-225}).
\item For every other entry, it checks that each declared file exists under
\code{data/<aoi>/} (\code{.../manifest.py:226-235}).
\item It re-checks the \code{expected\_files} union
(\code{.../manifest.py:237-240}).
\end{enumerate}

It is a \emph{declared-files-exist} check: it does not know which datasets a
given consumer needs. An AOI whose manifest has no \code{dw\_class\_fractions}
entry at all passes it with an empty missing list, and fails only later, in
the consumer, for a different reason.

% ---------------------------------------------------------------------------
\section{The catalogue and the \code{Slot} contract}
\label{sec:slots}

\subsection{\code{Slot}}

A \code{Slot}
(\code{mal-core/src/mal\_core/download/catalog.py:97-108}) is a frozen
dataclass with three fields: \code{key} (the manifest key),
\code{required} (whether absence is an error) and \code{per\_year} (whether
the entry is keyed by year, in which case the resolved value is a
\code{dict[year, Path]}). A consumer declares a mapping
\code{\{slot\_name: Slot(...)\}} and receives
\code{\{slot\_name: Path | dict[year, Path] | None\}}. The consumer's only
knowledge of the datasets is the manifest key.

\subsection{Resolution semantics}

\code{resolve\_inputs}
(\code{mal-core/src/mal\_core/download/catalog.py:176-268}) is the entry
point; \code{resolve} (\code{.../catalog.py:155-173}) is the single-key
version. Both delegate the per-entry lookup to
\code{resolve\_dataset\_file}
(\code{.../manifest.py:39-114}). There are exactly three failure modes,
distinguishable by their messages:

\begin{description}
\item[No manifest entry] An optional slot resolves to \code{None} --- the
signal for ``this input is absent, carry on without it''. A required slot
raises \code{FileNotFoundError} naming the AOI and telling the user to run
\code{malariasim download} (\code{.../catalog.py:228-235},
\code{...:258-263}).
\item[Entry present, file missing on disk] \code{FileNotFoundError}: a stale
manifest is an error, never a silent fallback (\code{.../manifest.py:103-114}).
The \code{strict} flag exists for diagnostics, but every consumer uses the
strict default.
\item[Per-year entry lacking the requested year] \code{FileNotFoundError}:
silently substituting another year's file would be a data-integrity bug, and
the message says so (\code{.../manifest.py:91-98}).
\end{description}

An explicit override always wins over the manifest and is validated to exist
(\code{.../catalog.py:207-215}).

\subsection{There is no conventional-name fallback}

The catalogue's docstring is unambiguous: there is no conventional-name
fallback inside consumers (\code{.../catalog.py:26-29}), and the rule is
restated at the point of resolution, in the docstring of
\code{resolve\_dataset\_file}
(\code{.../manifest.py:58-68}): a dataset that was never registered does not
exist as far as the pipeline is concerned. The only place conventional names are
derived is the migration path \code{register\_existing}
(\code{.../catalog.py:334-427}), which scans a pre-manifest data directory,
derives the conventional names from each loader's own \code{formats} plus the
runner's naming rules, registers what it finds, and never overwrites an entry
that already exists (\code{.../catalog.py:359-377}). It can be run as a dry
run (\code{dry\_run=True}), and it reports \code{registered}, \code{skipped},
\code{missing} and \code{dry\_run}
(\code{.../catalog.py:430-435}).

The migration path does embed a second, hand-written copy of the artefact
naming rules for the ingest products
(\code{.../catalog.py:411-421}), which are not derived from the producers
themselves. This is a known duplication and is listed in
Section~\ref{sec:discrepancies}.

\subsection{What \code{malariasim datasets} reports}

\code{catalog\_status} (\code{.../catalog.py:292-325}) produces one row per
known key: one row per downloaded dataset (derived from the registry) and one
row per ingest artefact (from \code{INGEST\_ARTIFACTS}). Each row carries the
key, the source that declares it, whether it is required for the ABM, whether
it is in the manifest, whether the declared file exists on disk, and the path.

The row shows only the \emph{first} declared file of an entry
(\code{.../catalog.py:305-318}). A multi-year product is therefore not fully
described by its row, and the \code{--missing} filter
(\code{mal-core/src/mal\_core/cli.py:830-831}) hides a partially complete
entry.

% ---------------------------------------------------------------------------
\section{Order of operations}
\label{sec:order}

\subsection{Pseudocode}

\begin{lstlisting}[basicstyle=\small,frame=single,breaklines=true]
run_download(aoi, datasets, outputs, years, months, output_dir, cache_dir):
    aoi_obj = AOI.from_slug(aoi);  registry = discover_downloaders()
    selected = registry restricted by --datasets, or all of registry
    out_dir = output_dir or data/<aoi>;  create it
    cache = cache_dir or None
    profile_mode = (datasets is None and outputs is None)
    for each (name, spec) in selected:
        if not credential_present(spec): record(name, skipped); continue
        try:
            for each (output_name, func) in spec.outputs:
                if outputs given and output_name not in outputs: continue
                if profile_mode and spec.abm_default_outputs is not None
                        and output_name not in spec.abm_default_outputs: continue
                if spec.is_time_series:
                    if no years given: warn and continue
                    format = spec.formats[output_name] or "monthly"
                    result = func(aoi=aoi_obj, years=years, months=months, cache_dir=cache)
                    if result is None: continue          # nothing to write
                    if format == "daily":        write ONE multi-year NetCDF
                    elif format == "monthly_nc": write ONE multi-year monthly NetCDF
                    else: for each year: reduce over time; write GeoTIFF
                else:
                    result = func(aoi=aoi_obj, cache_dir=cache)
                    if result is None: continue
                    write GeoTIFF or NetCDF
                register what was written
            record(name, ok)
        except Exception as e:
            log it;  record(name, error, str(e))          # next loader continues
    warn_manifest_limits(aoi)                  # advisory; never fails the run
    return the per-dataset status map
\end{lstlisting}

\subsection{Step, actor, order and mode}

The CLI parses the comma-separated lists first
(\code{cli.py:177-189}). The runner then builds the AOI object
(\code{runner.py:82-83}); the registry scans the loaders package, imports
each module and reads its \code{DOWNLOADER}
(\code{registry.py:47-107}); the runner selects loaders by name and creates
the output directory (\code{runner.py:87-93}). Per loader, the runner
probes each declared credential family (\code{runner.py:101-103},
\code{...:20-46}). Per output, it filters by \code{--outputs} and then by
the ABM profile, and chooses the write path from the declared format
(\code{runner.py:106-110}, \code{...:121-124}). The loader fetches and
reduces and returns an in-memory array (\code{runner.py:128}, \code{...:165},
\code{...:194}, \code{...:244}); the writer writes the file
(\code{writer.py:12-53}); the manifest is rewritten to register it
(\code{runner.py:147-159} and siblings). Last, the licences are aggregated and
any usage limits warned about (\code{runner.py:270-276}).

\subsection{Consequences of that order}

These are the falsifiable claims the order produces.

\begin{enumerate}
\item \textbf{Registration is per output, not per dataset.} Each output is
written and registered before the next output of the same loader is touched
(\code{runner.py:106-262}); a loader that fails on its third output leaves its
first two registered.
\item \textbf{A missing credential removes the whole loader, not one output.}
The gate runs before the output loop (\code{runner.py:101-103}), so the ERA5
loader, whose two outputs share the \code{cds} credential, contributes
neither or both.
\item \textbf{A loader that returns \code{None} writes nothing and registers
nothing} (\code{runner.py:129-130}, \code{...:166-167}, \code{...:195-196},
\code{...:246-247}).
\item \textbf{A failure is contained to its loader.} The exception is caught
at the loader level and reported as a status (\code{runner.py:266-268}); the
result is a per-dataset status map, not a pass/fail verdict for the AOI.
\item \textbf{The manifest is consistent after an interruption.} Because the
manifest is rewritten after every output, a killed run leaves a manifest
whose every declared file exists --- possibly an incomplete data set, but
never a manifest pointing at a missing file.
\item \textbf{Nothing is skipped because it already exists.} The runner does
not consult the manifest before downloading; idempotence comes from the
loaders' own caches.
\item \textbf{The output directory is created before any credential check}
(\code{runner.py:92-93}), so a run that skips everything still leaves an
empty \code{data/<aoi>/}.
\end{enumerate}

\subsection{A described situation}

\emph{A reader asks for the daily rainfall for 2024 over the Ghana AOI.} The
reader runs \code{malariasim download --aoi ghana --datasets chirps --years
2024}. The CLI splits the comma-separated strings
(\code{cli.py:177-189}) and calls the runner. The runner builds the AOI
object for \code{ghana} and scans the loaders package; the registry returns
every loader, and the runner keeps only \code{chirps}. The CHIRPS loader
declares no credential, so the gate passes. Its single output
\code{rainfall\_daily} survives the filters, and because \code{--datasets}
was named, the profile mode is off. The loader is time series, so years were
required and were given; its declared format is \code{"daily"}
(\code{.../chirps.py:346-348}), so the multi-year NetCDF branch is taken. The
loader walks the requested days, fetching one public GeoTIFF per day into its
cache, reprojecting each onto the AOI grid, and returns one array with
dimensions \code{(time, y, x)}. The runner derives the NetCDF variable name
from the output name with the \code{\_daily} suffix stripped --- \code{rainfall}
--- because the returned array carries no name of its own
(\code{runner.py:135-139}). The writer writes
\code{data/ghana/ghana\_rainfall\_daily\_2024\_2024\_daily.nc}
(\code{runner.py:140-142}) and the manifest gains an entry keyed
\code{chirps\_rainfall\_daily} with the period of the returned time axis
(\code{runner.py:147-159}). The CLI prints one status line and the
per-dataset status map.

\section{Authentication and credential behaviour}
\label{sec:auth}

\subsection{The gate}

\code{\_check\_auth} (\code{mal-core/src/mal\_core/download/runner.py:20-46})
probes one artefact per credential family:

\begin{longtable}{@{}llp{0.44\textwidth}@{}}
\toprule
Family & Probe & Missing-credential behaviour \\
\midrule
\endhead
\code{cds} & the file \code{\textasciitilde/.cdsapirc}, or the path in
\code{CDSAPI\_RC} (\code{runner.py:22-33}) & warning, then the loader is
skipped. The probe catches \code{OSError} as well as a missing file, so a
sandboxed home directory reads as "no auth". \\
\code{earthdata} & the environment variable \code{EARTHDATA\_TOKEN} is set
(\code{runner.py:34-37}) & warning, then the loader is skipped. \\
\code{earthengine} & the file \code{\textasciitilde/.config/earthengine/credentials} exists
(\code{runner.py:38-45}) & warning naming the \code{earthengine authenticate}
command, then the loader is skipped. \\
\code{none} & nothing (\code{runner.py:21-46}) & never skipped. \\
\bottomrule
\end{longtable}

A missing credential never raises. The loader is skipped, the status is
\code{\{"status": "skipped", "reason": "auth missing"\}}
(\code{runner.py:102}), and the run continues with the next loader. A user
without any credentials still receives a report; the report says which
datasets were skipped and why.

\subsection{Inside a directly-called loader}

The gate is a convenience of the runner, not a property of the loaders. A
loader invoked directly raises instead of skipping: the Earth Engine loader
raises \code{RuntimeError} naming the project and the one-line fix
\code{earthengine authenticate}
(\code{mal-commonlib/src/mal\_commonlib/data/loaders/dynamic\_world.py:149-159}),
never opening an interactive browser flow; the ERA5 loader raises
\code{RuntimeError} when the CDS client cannot be constructed, naming the
expected configuration (\code{.../era5.py:201-208}, and again for the wind
product at \code{...:402-409}).

The Earth Engine loader also installs a shim around the OAuth module's
application-default-credentials probe, so that an unreadable Google Cloud
credentials path cannot abort initialisation when the Earth Engine token
itself is fine (\code{.../dynamic\_world.py:129-147}). A filesystem permission
problem is therefore reported as an initialisation failure, not as a missing
credential.

\subsection{The credential artefact, per provider}

\begin{longtable}{@{}lll@{}}
\toprule
Provider family & Artefact that must exist & Used by \\
\midrule
\endhead
Copernicus CDS & \code{\textasciitilde/.cdsapirc} (or \code{CDSAPI\_RC}) & ERA5 (air
temperature and wind) \\
NASA Earthdata & \code{EARTHDATA\_TOKEN} in the environment & MODIS NDVI, SMAP
salinity \\
Google Earth Engine & \code{\textasciitilde/.config/earthengine/credentials} & Dynamic World
fractions \\
none & --- & CHIRPS, JRC GSW, WorldPop, GLW, GHSL, Overture buildings, the
wildlife proxy, GSHHG coastline \\
\bottomrule
\end{longtable}

The Earth Engine project identifier used by the loader is a module-level
constant in the loader file itself, not a value from the constants file
(\code{.../dynamic\_world.py:57}).

% ---------------------------------------------------------------------------
\section{The write path}

\subsection{Three write paths}

The declared format of an output selects the path
(\code{mal-core/src/mal\_core/download/runner.py:121-124}; the default is
\code{"monthly"}).

\begin{itemize}
\item \code{"daily"}: one multi-year NetCDF covering the whole requested year
span, named with the first and last requested years (\code{runner.py:140}).
The period is read from the returned array's time coordinate --- \code{time}
if present, otherwise \code{valid\_time} (\code{runner.py:144-146}) --- and
recorded in the manifest (\code{runner.py:147-159}).
\item \code{"monthly\_nc"}: the same with a \code{\_monthly.nc} name and
monthly steps; used by the SMAP salinity product (\code{runner.py:161-191},
\code{.../smap.py:433-435}).
\item Anything else: per-year GeoTIFFs. If the returned array has a time
dimension, the runner slices it per year and reduces over time before
writing; if it has none, it writes the single array under the first requested
year (\code{runner.py:220-262}).
\end{itemize}

Static products take their extension from the returned object: a NetCDF for
an \code{xarray.Dataset}, a GeoTIFF for anything else
(\code{runner.py:249-250}).

\subsection{The writer}

\code{save\_product} (\code{mal-core/src/mal\_core/download/writer.py:12-53})
is the single save point. It infers the format when none is given
(\code{.../writer.py:29-30}), writes GeoTIFFs through the \code{rio} accessor
and NetCDFs through \code{xarray}
(\code{.../writer.py:32-49}), and refuses a three-dimensional array sent to
the GeoTIFF writer with a \code{ValueError}
(\code{.../writer.py:32-40}). Loaders never write files: they return data in
memory, and the runner passes it to the writer.

\subsection{The NetCDF variable name is a contract}

The name of the variable inside a NetCDF file is how the ingest stage finds
the channel. The runner prefers the loader's own array name and falls back to
the output name with the \code{\_daily} suffix stripped
(\code{runner.py:135-139}); the monthly path strips \code{\_salinity}
instead (\code{runner.py:170}). The writer applies the name when it turns an
array into a dataset (\code{writer.py:45-48}).

In practice: ERA5's temperature loader renames its array to
\code{air\_temp\_c}
(\code{mal-commonlib/src/mal\_commonlib/data/loaders/era5.py:335},
\code{...:337}), so the file carries \code{air\_temp\_c}; CHIRPS returns an
unnamed array, so the fallback yields \code{rainfall}, which is the channel
the ingest builder reads
(\code{mal-core/src/mal\_core/ingest/daily\_nc.py:198-199}); the JRC loader
returns a dataset whose variable is \code{jrc\_persistent}
(\code{.../jrc\_gsw.py:318-321}).

% ---------------------------------------------------------------------------
\section{The AOI reduction: where the aggregation actually happens}
\label{sec:reduction}

\subsection{The seam}

The runner never resamples. It calls the loader, and the loader returns
something already on the AOI grid. The ``areal fraction before reduction''
rule is therefore a \emph{per-loader implementation}, not a framework
guarantee --- and neither is the grid. Each loader derives the grid for
itself by calling \code{aoi.cells\_per\_side()}
(\code{mal-commonlib/src/mal\_commonlib/aoi.py:195-199}); neither the runner
nor the writer checks the shape or the CRS of the array it is handed. Both
the grid and the rule are per-loader conventions: a new loader could
satisfy the grid and violate the rule, or the other way round, and nothing
in the framework would notice.

\subsection{The Dynamic World rule, stated exactly}

The Dynamic World loader is the reference implementation of the rule
(\code{mal-commonlib/src/mal\_commonlib/data/loaders/dynamic\_world.py}):

\begin{enumerate}
\item For each run year, the Earth Engine image collection
\code{GOOGLE/DYNAMICWORLD/V1} is filtered to that year and its \code{label}
band is reduced by mode over the year's Sentinel-2 passes
(\code{.../dynamic\_world.py:235-246}). Cloudy passes contribute transient
labels that lose the vote.
\item The modal-label raster is requested at a $100\,\mathrm{m}$ scale, not
the native $10\,\mathrm{m}$: the request is a hundred times cheaper and the
$1\,\mathrm{km}$ aggregation smooths the difference anyway
(\code{.../dynamic\_world.py:56-63}).
\item Each label code is turned into a binary layer and reprojected onto the
AOI grid with an \emph{averaging} resampling, so each cell receives the
fraction of its area covered by that class
(\code{.../dynamic\_world.py:302}, \code{...:314}).
\item The class fractions are summed into four feature columns and divided by
the number of valid pixels in the cell
(\code{.../dynamic\_world.py:318-331}). A cell with no valid pixel gets the
NoData sentinel \mskey{ingest.nodata_sentinel}; the divisor is floored at
\mskey{ingest.valid_count_epsilon} so that the division cannot divide by
zero.
\item The four feature columns are \code{lc\_built},
\code{lc\_wetland}, \code{lc\_crop} and \code{lc\_open\_plus\_water}, the last
being the union of grass, shrub and scrub, bare and water
(\code{.../dynamic\_world.py:92-103}).
\end{enumerate}

What this reduction loses: the sub-$100\,\mathrm{m}$ composition of every
cell; the identity of the minority class in a mixed cell; and the difference
between a class that dominates because it is persistent and one that won a
close annual vote --- not recoverable downstream, because the annual mode has
already discarded the vote distribution.

The loader returns one frame per requested year on a \code{time} dimension
(\code{.../dynamic\_world.py:378-410}); the ingest stage averages the frames
into the single static plane of the habitat contract.

\subsection{The JRC rule, stated exactly}

The JRC loader streams the \code{occurrence} asset from the Planetary
Computer STAC catalogue and bins it at a fixed occurrence threshold
(\code{.../jrc\_gsw.py:61-72}) --- the product's own definition of
"permanent water". The binarised mask and the valid-pixel mask are each
reprojected onto the AOI grid and their ratio is the cell's permanent-water
fraction (\code{.../jrc\_gsw.py:209-235}). The year used is the last year of
the product's coverage range, and the loader refuses a year outside it
(\code{.../jrc\_gsw.py:296-299}).

\subsection{Resampling, per product class}

\begin{longtable}{@{}llp{0.48\textwidth}@{}}
\toprule
Product class & Resampling & Consequence \\
\midrule
\endhead
Continuous field (ERA5 temperature) & bilinear
(\code{.../era5.py:151}) & smooths across the coast; before the warp the
loader fills the NaN holes with the \emph{nearest finite neighbour} (a
distance transform), because GDAL's bilinear warper turns NaN into an
expanding nodata fillet; only afterwards is the NoData sentinel written,
as the raster's own nodata tag rather than as a substitute value
(\code{.../era5.py:137-146}) \\
Rainfall (CHIRPS) & nearest (\code{.../chirps.py:230-241}) & the loader's
comment states the reason: for small coastal AOIs the source window is tiny
and mostly ocean NoData, and a bilinear blend propagates that NoData into
every destination cell \\
Categorical settlement class (GHSL) & nearest
(\code{.../ghsl.py:199-206}) & a categorical field must not be interpolated;
the class codes are preserved \\
Population (WorldPop) & sum-preserving rectangle sum
(\code{.../worldpop.py:141-200}) & the per-cell value stays a \emph{count of
people}, not a density; the loader's docstring records that an earlier
nearest-neighbour implementation discarded most of the source mass
(\code{.../worldpop.py:30-36}) \\
Livestock density (GLW) & nearest (\code{.../glw.py:183-189}) & density is
intensive; summing it would inflate the value by the number of fine cells per
coarse pixel \\
Land-cover class fractions (Dynamic World) & averaging
(\code{.../dynamic\_world.py:302}) & the value is an areal fraction by
construction \\
Building fraction (Overture) & centroid containment
(\code{.../buildings.py:18-25}) & the value is 1.0 in any cell containing at
least one footprint centroid; it is \emph{not} an area fraction \\
Coastline land fraction (GSHHG) & sub-cell rasterisation then averaging
(\code{.../coastline.py:33-37}) & the fraction of each cell that is land; no
dilation is applied \\
\bottomrule
\end{longtable}
\label{tab:resampling}

\subsection{What is lost in the population layer}

The WorldPop reduction preserves the total by summing source pixels into
destination cells, so the AOI-grid raster is a count per cell and the sum
over cells is the population of the covered area. What is lost is the
within-cell distribution: a cell's inhabitants are placed uniformly across it,
and a cell straddling an urban boundary keeps only its total. The source's
own resolution is finer than the target, so the aggregation is not lossy in
mass --- it is lossy in position.

% ---------------------------------------------------------------------------
\section{What arrives from outside, and what is decided here}

\begin{itemize}
\item \textbf{Which datasets exist} is decided by the loaders package, by
existence (\code{registry.py:12-25}).
\item \textbf{Which outputs a default run fetches} is decided by the
registry's profile table and the loaders' \code{abm\_default\_outputs}
(\code{registry.py:55-75}).
\item \textbf{Which years and months} are decided by the caller
(\code{cli.py:158-163}).
\item \textbf{Which AOI} is decided by the caller and resolved against
\code{aois.yaml}; \textbf{the grid resolution} is decided by the AOI
registry.
\item \textbf{The product's native resolution} is a fact about the upstream
provider and is not a choice here.
\item \textbf{The class composition and the areal-fraction rule} are decided
by the loader's source code (\code{dynamic\_world.py:92-103}).
\item \textbf{The occurrence threshold for ``permanent water''} is a module
constant in the JRC loader (\code{jrc\_gsw.py:68-72}).
\item \textbf{The year-aggregation rule for monthly products} is a module
constant defaulting to the mean (\code{runner.py:203-209}).
\item \textbf{The credential families} are decided by each loader's
\code{requires\_auth}.
\item \textbf{The licence and attribution lines} are decided by each loader
(\code{licenses.py:1-12}).
\item \textbf{The NoData sentinel} is not a choice here: it is shared with the
engine and read from the constants file
(\code{dynamic\_world.py:110-112}).
\end{itemize}

A reader who wants a different scenario --- a different year, a different
AOI, a different land-cover decomposition --- changes the corresponding row
in that table and nothing else. A reader who wants a different \emph{science}
changes the constants file or the loader. The stage does have parameters of
its own, though: the credential environment variables \code{CDSAPI\_RC} and
\code{EARTHDATA\_TOKEN} (\code{runner.py:26}, \code{...:35}), the request
padding variable \code{ERA5\_AREA\_PAD\_DEG}
(\code{.../era5.py:176}), the cache root \code{XDG\_CACHE\_HOME}
(\code{.../chirps.py:59}), and the CLI's \code{--years}, \code{--months} and
\code{--output-dir} (\code{mal-core/src/mal\_core/cli.py:161-163}).

% ---------------------------------------------------------------------------
\section{Provenance of the values}
\label{sec:provenance}

\subsection{Three tiers, and where each class of value lives}

\begin{itemize}
\item \textbf{Model parameters shared with the engine} --- the NoData
sentinel (\mskey{ingest.nodata_sentinel}), its integer twin
(\mskey{ingest.nodata_sentinel_int}), the valid-count epsilon
(\mskey{ingest.valid_count_epsilon}), the wildlife suitability table
(\mskey{wildlife.class_suitability.trees} and its eight siblings) and the
wildlife component weights (\mskey{wildlife.weights.habitat},
\mskey{wildlife.weights.water}, \mskey{wildlife.weights.remoteness}) ---
are a \emph{model definition}. They live in
\code{mal-core/config/model\_constants.yaml} and are read through the
read-only accessor; the loaders read them and never define their own --- both
the Dynamic World loader
(\code{.../dynamic\_world.py:112}) and the JRC loader
(\code{.../jrc\_gsw.py:83}) read the valid-count epsilon from the same
configuration key.
\item \textbf{Upstream facts} --- native resolutions, product identifiers,
coverage years, class codes, the occurrence threshold, the request scale, the
request padding --- are \emph{measured or literature-anchored}. They are
stated by the provider and recorded in the loader's module constants and
docstrings; they are cited to the source in this document and are not
generated.
\item \textbf{Engineering choices} --- the ABM profile table, the default
year-aggregation rule, the file-naming patterns, the cache locations --- are
an \emph{engineering prior}: chosen by the project, revisable, and cited to
the module that holds them.
\end{itemize}

The following values in this document come from the generated macros, that
is, from \code{mal-core/config/model\_constants.yaml} through
\code{docs/concepts/generated/macros.tex}: the NoData sentinels
\mskey{ingest.nodata_sentinel} and \mskey{ingest.nodata_sentinel_int};
the valid-count epsilon \mskey{ingest.valid_count_epsilon}; the wildlife
suitability table \mskey{wildlife.class_suitability.water},
\mskey{wildlife.class_suitability.trees},
\mskey{wildlife.class_suitability.grass},
\mskey{wildlife.class_suitability.flooded_vegetation},
\mskey{wildlife.class_suitability.crops},
\mskey{wildlife.class_suitability.shrub_and_scrub},
\mskey{wildlife.class_suitability.built},
\mskey{wildlife.class_suitability.bare} and
\mskey{wildlife.class_suitability.snow_and_ice}; the wildlife component
weights \mskey{wildlife.weights.habitat}, \mskey{wildlife.weights.water} and
\mskey{wildlife.weights.remoteness}; the ingest-side coastline parameters
\mskey{ingest.coastline_buffer_m} and \mskey{ingest.coastline_resolution};
the cell-size hint \mskey{ingest.cell_size_hint_m}; and the contract
versions \mskey{contracts.env_nc} and \mskey{contracts.cogs_state}.

Every other magnitude quoted in this document is a \emph{reading of the
code}: a number typed into a loader's source or into a docstring, cited to
that file and line. Such readings are explicitly marked as such in the
evidence appendix.

\subsection{The output contracts of the pipeline}

The generated contract table below is authoritative for the pipeline's
artefacts; the download stage participates in it through the \code{manifest}
row.

% Generated by tools/gen_concepts_values.py -- do not edit.
\paragraph{The output contracts}
\begin{longtable}{@{}llllll@{}}
\toprule
Artefact & Producer & Consumers & Version & Bands / variables & Sentinel \\
\midrule
\endhead
\code{env} & \code{mal\_core.ingest.env.build\_env\_tensor} & abm, prediction, scoring & "3.0" & rainfall, air\_temp\_c, jrc\_persistent, lc\_built, lc\_wetland, lc\_open\_plus\_water, lc\_crop, land\_frac (optional), salinity\_ppt (optional) & none \\
\code{habitat} & \code{mal\_core.ingest.env.build\_env\_tensor} & abm &  & hab\_type, habitat\_potential, row, col & none \\
\code{state\_cog} & \code{mal\_abm\_fast::write\_state\_cog} & scoring, prediction, training & "4.0" & adult\_patch\_count, adult\_cell\_count & -9999.0 \\
\code{state\_sidecar} & \code{mal\_abm\_fast::write\_state\_sidecar} & scoring & "4.0" &  &  \\
\code{manifest} & \code{mal\_core.download.manifest.update\_dataset} & download, ingest, abm &  &  &  \\
\code{host\_static} & \code{mal\_core.ingest.hosts.build\_host\_dataset} & abm &  & humans, livestock, urban, buildings &  \\
\code{mobility\_day} & \code{mal\_core.ingest.mobility.build\_mobility\_dataset} & abm &  &  &  \\
\bottomrule
\end{longtable}

\paragraph{The state sidecar's keys, in emission order}
\begin{longtable}{@{}l@{}}
\toprule
Key \\
\midrule
\endhead
\code{crs} \\
\code{transform} \\
\code{aoi\_slug} \\
\code{scale} \\
\code{year} \\
\code{month} \\
\code{seed} \\
\code{n\_rollouts} \\
\code{rollout\_index} \\
\code{generator\_version} \\
\code{abm\_params\_hash} \\
\code{contract\_version} \\
\code{band\_names} \\
\code{nodata} \\
\code{shape} \\
\bottomrule
\end{longtable}


% ---------------------------------------------------------------------------
\section{Outputs and their consumption contract}
\label{sec:consumption}

\subsection{What the download stage itself produces}

Four kinds of file, all under \code{data/<aoi>/}, all registered in the
manifest:

\begin{itemize}
\item A multi-year daily NetCDF, \code{<aoi>\_<output>\_<start>\_<end>\_daily.nc},
registered as \code{"time-series"} with \code{format="nc"} and a
\code{period}.
\item A multi-year monthly NetCDF,
\code{<aoi>\_<output>\_<start>\_<end>\_monthly.nc}, registered the same way.
\item A per-year GeoTIFF, \code{<aoi>\_<output>\_<year>.tif}, registered as
\code{"time-series"} with \code{format="tif"} and keyed by year.
\item A static GeoTIFF or NetCDF, \code{<aoi>\_<output>.<ext>}, registered as
\code{"static"}.
\end{itemize}

The consumption contract for all of them is the same and is enforced by the
catalogue: a consumer asks for a manifest key, and the manifest answers with
a path; if the path is absent the consumer is told which dataset is missing.

\subsection{What the ingest stage requires}

The ingest stage declares its inputs as \code{Slot}s and nothing else
(\code{mal-core/src/mal\_core/ingest/daily\_nc.py:87-98}):

\begin{longtable}{@{}lll@{}}
\toprule
Slot & Manifest key & Required? \\
\midrule
\endhead
\code{rainfall} & \code{chirps\_rainfall\_daily} & yes \\
\code{air\_temp} & \code{era5\_air\_temp\_daily} & yes \\
\code{jrc\_features} & \code{jrc\_habitat\_features} & yes \\
\code{wc\_fractions} & \code{dw\_class\_fractions} & yes \\
\code{land\_mask} & \code{coastline\_land\_mask} & no \\
\code{host\_static} & \code{host\_static} & no \\
\code{salinity} & \code{smap\_salinity} & no \\
\bottomrule
\end{longtable}

A missing required slot is a hard error at resolution time
(\code{.../daily\_nc.py:113-122}); a missing optional slot disables the
corresponding behaviour. Note what this means: the ingest stage re-reads the
static rasters at the target shape, resampling bilinearly if the grid differs
(\code{.../daily\_nc.py:126-148}), so the download stage's grid is a
convention for those products rather than an enforced contract.

\subsection{What the ABM requires}

The ABM wrapper validates completeness first
(\code{mal-core/src/mal\_core/abm/wrapper.py:116-121}) and then resolves
seven paths: \code{env}, \code{habitat}, \code{host\_static}, \code{wind},
\code{mobility\_day}, \code{mobility\_night} and \code{livestock\_mobility}
(\code{...:141-165}). Only \code{env} and \code{habitat} are mandatory; the
others are passed to the binary when present
(\code{...:167-178}). The wind path is resolved by the literal key
\code{"wind"} rather than through a loader declaration, and the source says
so in a comment (\code{...:149-152}). This is the one consumer that
hard-codes a dataset key.

The \code{required\_for\_abm} flag is \emph{not} the gate the ABM uses. It is
declared only by the Dynamic World and ERA5 loaders
(\code{.../dynamic\_world.py:500}, \code{.../era5.py:506}), copied into the
manifest entry, and displayed by \code{malariasim datasets}; the gate itself
is completeness plus the presence of \code{env} and \code{habitat}.

\subsection{Ordering constraint with the rest of the pipeline}

The download stage must precede ingest, and ingest must precede the ABM. The
\code{abm} command chains the first two automatically: it calls the download
runner for the datasets the contract consumes and then builds the env NetCDF
(\code{mal-core/src/mal\_core/cli.py:665-693}). The manifest carries raw
layers and derived products alike --- the \code{env} and \code{habitat} keys
are registered by ingest into the same file the download stage writes
(\code{mal-core/src/mal\_core/download/catalog.py:55-68}).

% ---------------------------------------------------------------------------
\section{Cache, state and restart}
\label{sec:state}

\subsection{Three kinds of on-disk thing}

\begin{description}
\item[Cache] \code{\textasciitilde/.cache/mal\_commonlib/<dataset>/} (or
\code{\$XDG\_CACHE\_HOME/mal\_commonlib/<dataset>/}), overridable per loader
(\code{.../era5.py:69-72}, \code{.../chirps.py:57-61},
\code{.../worldpop.py:105-107}, \code{.../glw.py:97-99},
\code{.../buildings.py:52-54}, \code{.../dynamic\_world.py:170-176}). It
holds the network-fetched artefacts --- per-month NetCDFs, per-day GeoTIFFs,
Dynamic World label tiles, rasterised \code{.npz} results. It is a pure
cache: deleting it changes nothing about the manifest.
\item[Products] \code{data/<aoi>/}: the AOI-grid reductions, named by the
conventions above and registered in the manifest.
\item[Manifest] \code{data/<aoi>/manifest.json}: the state.
\end{description}

The runner passes a \code{cache\_dir} argument to every loader; when it is
not given, each loader falls back to its own location
(\code{mal-core/src/mal\_core/download/runner.py:95},
\code{...:128}, \code{...:244}). A cached file makes a re-run cheap; a cold
cache makes it slow, but the products and the manifest are the same either
way.

\subsection{The restart test}

If the stage were interrupted and restarted, what would have to be on disk to
resume identically? The answer is: the manifest, plus the files it names.

Because the manifest is rewritten after every single output
(\code{runner.py:147-159} and siblings) and \code{expected\_files} is
recomputed on every write (\code{.../manifest.py:172-180}), an interrupted
run leaves a manifest that is internally consistent: every file it declares
exists, and the set of files may be incomplete. Re-running the same command
re-fetches what is missing. It does not \emph{skip} what is present --- the
runner always re-downloads, and the loaders' caches are what make that cheap.

The one thing that is \emph{not} restart-safe is a half-written product file:
the write and the manifest update are separate steps, so a crash between them
leaves an orphan file that no manifest entry names. The next run overwrites
it.

% ---------------------------------------------------------------------------
\section{Licensing and attribution}

Each loader declares a licence identifier and a literal attribution line
(\code{mal-commonlib/src/mal\_commonlib/data/loaders/era5.py:481-491},
\code{.../chirps.py:335-340}, and the equivalent in every other loader).
The identifiers come from a registry of known licences with their usage
limits (\code{mal-commonlib/src/mal\_commonlib/data/licenses.py:25-146}).

The runner copies both into the manifest entry
(\code{runner.py:147-159}, \code{...:212-223}) and, at the end of the run,
aggregates the licences of all entries and warns about any usage limits
(\code{runner.py:270-276}). The warning is advisory: it is wrapped in a
\code{try}/\code{except} whose comment states that a download must never
break over a warning. Derived products registered by ingest carry compound
identifiers joined with \code{+} (\code{.../licenses.py:110-124}).

The attribution stored for the wildlife proxy still names ESA WorldCover as a
source, although the loader now reads the Dynamic World modal-label raster
(\code{.../wildlife.py:277-285}, \code{...:wildlife.py:42-45}). The credit
line therefore does not match the data. This is recorded in
Section~\ref{sec:discrepancies}.

% ---------------------------------------------------------------------------
\section{What this stage does not do}
\label{sec:absences}

These are absences of implementation, not of documentation.

\begin{enumerate}
\item \textbf{No transformation into model inputs.} The env tensor and the
habitat GeoPackage are ingest products
(\code{mal-core/src/mal\_core/download/catalog.py:55-68}).
\item \textbf{No resampling for the consumer.} The runner writes what the
loader returned; the loaders own the reduction
(Section~\ref{sec:reduction}).
\item \textbf{No validation of scientific content.} Only existence and
declaration are checked (\code{.../manifest.py:192-242}).
\item \textbf{No skipping of already-downloaded outputs.} The runner always
re-fetches and re-registers; idempotence comes from the caches.
\item \textbf{No retry or back-off policy at the runner level.} A loader that
raises is recorded as an error and the run moves on
(\code{runner.py:266-268}).
\item \textbf{No load-or-download fast path.} There is no function that
checks the manifest and returns an existing file instead of fetching. The
loaders always fetch; the manifest is written by the runner and read by
consumers. The convention described in the project's \code{AGENTS.md} as
``load-or-download'' is not implemented in the code.
\item \textbf{No \code{--all} flag.} The two places that advertise it
(\code{cli.py:171-173}, \code{abm/wrapper.py:118-121}) refer to a flag that
does not exist.
\item \textbf{No configurable year-aggregation rule in practice.} The runner
supports \code{"sum"}, \code{"first"} and the default \code{"mean"}
(\code{runner.py:203-209}), but no loader declares a \code{reductions} map,
so the alternative branches are unreachable and every monthly product is
reduced by the mean.
\item \textbf{No consumer-side dataset knowledge beyond manifest keys.} The
one exception is the ABM wrapper's literal \code{"wind"}
(\code{abm/wrapper.py:149-152}).
\item \textbf{No verification that the products agree with each other.} Two
products describing the same AOI are reduced independently and never
cross-checked.
\end{enumerate}

% ---------------------------------------------------------------------------
\section{Evidence appendix}

\subsection{Implementation notes a replicator needs}

\begin{enumerate}
\item \textbf{The runner is the only writer.} Loaders return arrays; the writer
is the only place that touches the filesystem for products
(\code{writer.py:12-53}). Write, then register, per output --- not per
dataset (\code{runner.py:147-159}).
\item \textbf{The output name, not the manifest key, names the file.} A
renamed loader changes the file names but not the manifest keys, so a
migration is then needed (\code{runner.py:49-66}).
\item \textbf{The NetCDF variable name is chosen by the runner} unless the
loader names its array (\code{runner.py:135-139}).
\item \textbf{The period is derived from the returned time coordinate.} A
loader with no time dimension cannot be registered with a period, and its
completeness check falls back to plain file existence
(\code{runner.py:144-146}, \code{.../manifest.py:205-225}).
\item \textbf{The \code{formats} map is a routing token, not a
description.} Dynamic World is one frame per year yet takes the ``daily''
path, because that path writes one multi-year NetCDF
(\code{registry.py:76-82}, \code{.../dynamic\_world.py:491}).
\item \textbf{The ERA5 wind output takes the ``daily'' path too}, although the
product is six-hourly: the loader returns a six-hourly dataset and the
runner writes it under a \code{\_daily} file name, so the stored file's
temporal resolution is the loader's, not the file name's
(\code{.../era5.py:493-496}, \code{...:era5.py:447-470}).
\item \textbf{The ERA5 request area is padded} by half a degree by default,
because a bare small bounding box produced a nearly empty source slice; the
padding is overridden by an environment variable, and a source slice less
than half finite is refused (\code{.../era5.py:56}, \code{...:158-186},
\code{...:120-127}).
\item \textbf{The ERA5 wind request area is hard-coded to a Ghana-wide
window}, not to the AOI's bounding box; the result is reprojected afterwards
(\code{.../era5.py:357}, \code{...:429}).
\item \textbf{The channel names} come from the loader where it names its
array: ERA5 returns \code{air\_temp\_c}, CHIRPS returns an unnamed array and
the runner supplies the fallback \code{rainfall}
(\code{.../era5.py:335}, \code{runner.py:135-139}).
\item \textbf{The wildlife proxy is the one loader that consumes other
products.} It reads the Dynamic World modal-label raster, the JRC occurrence
fraction and the Overture building fraction and combines them with the
weights \mskey{wildlife.weights.habitat},
\mskey{wildlife.weights.water} and \mskey{wildlife.weights.remoteness}. A
component that fails to load silently becomes a neutral constant --- $0.5$
for habitat and water, $1.0$ for remoteness (\code{.../wildlife.py:21-23},
\code{.../wildlife.py:200-233}).
\item \textbf{The JRC loader's year is bounded} by the product's coverage
range (\code{.../jrc\_gsw.py:74-76}, \code{...:289-293}); \textbf{the Dynamic
World loader refuses years before the record starts}
(\code{.../dynamic\_world.py:106}, \code{...:361-364}).
\item \textbf{Modis, SMAP and the retired hydrorivers loader} are the three
special cases of the profile table: SMAP and hydrorivers are entries with
empty lists, and \code{modis} is not an entry of the table at all, which is
why it is not excluded
(\code{registry.py:55-75}).
\end{enumerate}

\subsection{Readings of the code, not configured values}

The following counts and enumerations in this document are readings of the
source at the time of writing, not values from the constants file:

\begin{itemize}
\item The registry discovers twelve loader modules, listed in
Section~3.4; the retired loaders under
\code{mal\_commonlib/data/loaders/\_legacy/} are invisible to the scan
because they form a package (\code{registry.py:12-25}).
\item That table lists seventeen output-to-key mappings.
\item The profile table has twelve entries, of which ten assign outputs and
two are empty lists (\code{registry.py:55-75}).
\item The \code{INGEST\_ARTIFACTS} table holds eight derived artefacts
(\code{catalog.py:55-92}); the ingest stage declares four required and three
optional slots (\code{ingest/daily\_nc.py:87-98}); the ABM wrapper resolves
seven paths (\code{abm/wrapper.py:137-163}).
\end{itemize}


\subsection{Verified and not verified}

\begin{itemize}
\item \emph{Verified by reading the code}: the manifest schema and its three
top-level keys, checked against the real \code{data/ghana/manifest.json}
(\code{.../download/manifest.py:30-189}); the resolution semantics and their
three failure modes (\code{.../download/catalog.py:155-268},
\code{.../manifest.py:39-114}); the order of operations
(\code{.../download/runner.py:69-278}); the credential probes
(\code{.../download/runner.py:20-46}); the areal-fraction rule for Dynamic
World (\code{.../dynamic\_world.py:273-334}); the resampling choices in
Table~\ref{tab:resampling}.
\item \emph{Read but not executed}: the behaviour of the CDS client and of the
Earth Engine initialiser inside the loaders, and the Dynamic World
reduction path --- both were read, not run.
\item \emph{Not verified here}: the real numbers of a Ghana download (this
review did not run a download; the manifest on disk is the evidence used for
the naming and schema claims), and the scientific adequacy of the products,
which is not this stage's concern.
\end{itemize}

\subsection{Discrepancies found while writing this document}
\label{sec:discrepancies}

These are places where the code, the documentation and the generated tables
disagree. The document describes the code; the disagreements are listed here
so that they are not lost.

\begin{itemize}
\item \code{cli.py:171-173} and \code{abm/wrapper.py:118-121} advertise a
\code{--all} flag the CLI does not define; the behaviour is obtained by
omitting \code{--datasets} and \code{--outputs}.
\item \code{registry.py:55-75} and \code{runner.py:109-110}: \code{modis} is
absent from the profile table and declares no \code{abm\_default\_outputs},
so its value is \code{None} and the filter does not exclude it --- a default
run fetches MODIS NDVI, which no current consumer reads.
\item \code{runner.py:203-209}: the \code{"sum"} and \code{"first"} reduction
branches are unreachable, because no loader declares a \code{reductions} map.
\item \code{tools/gen\_concepts\_values.py:334}: the generated dataset
table's column headed ``Manifest keys'' enumerates the \emph{keys} of each
loader's \code{manifest\_keys} mapping, which are the output names; the
manifest keys are the mapping's \emph{values} (Table~\ref{tab:loaders}).
\item \code{.../wildlife.py:277-285}: the attribution still credits ESA
WorldCover, which is no longer the source.
\item \code{.../dynamic\_world.py:491} and \code{.../era5.py:493-496}: two
products are routed through the ``daily'' write path although one is annual
and the other six-hourly; the path chooses the container and the naming, not
the temporal resolution.
\item \code{.../era5.py:357}: the ERA5 wind request area is a hard-coded
Ghana-wide window rather than the AOI's bounding box.
\item \code{catalog.py:411-421}: the migration path embeds a hand-written
copy of the ingest artefacts' naming rules, not derived from the producers.
\item \code{abm/wrapper.py:149-152}: the ABM resolves the ERA5 wind by the
literal key \code{"wind"} instead of through the loader's declaration.
\item \code{.../manifest.py:91-98} against \code{.../manifest.py:103-114}: the
per-year miss and the stale-manifest miss both raise
\code{FileNotFoundError}; only their messages distinguish them.
\end{itemize}

\subsection{What this document must not be used for}

It is not an installation guide for the credential systems, not a licence
audit (that is \code{docs/licenses.md}), not a specification of the manifest
schema (which the code owns), and not a claim that the data in
\code{data/<aoi>/} is scientifically adequate for the malaria question. It
describes one stage: how external data arrives, is reduced onto the AOI grid,
is written, and is recorded.

\end{document}
